%%%%%%%%%%%%%%%%%%%%%%%%%%%%%%%%%%%%%%%%%%%%%%%%%%%%%%%%%%%%%%%%%%%%%%%
%                                                                     %
%                              Chapter 12                             %
%                                                                     %
%%%%%%%%%%%%%%%%%%%%%%%%%%%%%%%%%%%%%%%%%%%%%%%%%%%%%%%%%%%%%%%%%%%%%%%


\chapter{Case Studies}

To enable you to further explore the ways the concepts of calculus are
used as analytical tools in scientific and mathematical
investigations, this chapter presents four extended case studies.  The
four can be studied separately, although the first two and the last
two are loosely linked
\begin{description}

\item[Stirling's Formula] As an example of the way many of the
  ideas---Taylor series, numerical integration, reduction formulas,
  limits---developed in the earlier chapters of this book can be used
  in a tightly-reasoned argument to produce some powerful mathematical
  insights, in the first section we derive a famous formula
  approximating $n!$.  This formula is then applied to the binomial
  probability distribution.

\item[The Poisson Distribution] Chapter~12.2 continues the probability
  theme by developing the Poisson distribution and using it to study
  the frequency of radioactive decay events.

\item[The Power Spectrum] Chapter~12.3 builds on the study of
  periodicity begun in chapter 7. We develop the Fourier transform, a
  basic tool in the sciences for detecting the relative strength of
  periodic components in a noisy data set.

\item[Fourier Series] Chapter~12.4 expands on some of the ideas in
  chapter 11.  Here we develop tools for approximating functions over
  intervals using sums of sine and cosine terms.  This is an
  extensively used method in a wide range of disciplines, from
  thermodynamics to music synthesis.
\end{description} 
	
\newpage
%%%%%%%%%%%%%%%%%%%%%%%%%%%%%%%%%%%%%%%%%%%%%%%%%%%%%%%%%%%%%%%%%%%%%%%
%                                                                     %
%                              Section 1                              %
%                                                                     %
%%%%%%%%%%%%%%%%%%%%%%%%%%%%%%%%%%%%%%%%%%%%%%%%%%%%%%%%%%%%%%%%%%%%%%%

\section{Stirling's Formula}
\index{Stirling's formula}
\index{factorial}
\index{$n"!$}

Given a positive integer $n$, we define $n!$---pronounced {\em $n$
  factorial}---by the rule $n!=1\cdot 2\cdot 3\cdots (n-1)\cdot n$\,.
This is a convenient concept which occurs in a number of settings,
%
\mar{Factorials\\ in probability}	
%
particularly combinatorial and probabilistic ones.  For instance, the
probability\index{probability} of getting exactly $n$ heads out of
$2n$ tosses of a coin turns out to be
	$$
\dfrac{(2n)!}{2^{2n}(n!)^2}.
	$$ 
Unfortunately, evaluating $n!$ for values of $n$ at all large is
cumbersome at best.  Although though many calculators will compute
factorials, few of them can handle numbers as large as $1000!$\,.
Even when we can evaluate $n!$, we are often as interested in the
asymptotic behavior of a certain expression as much as in its exact
value for specific $n$.
%
\mar{$n!$ is difficult \\to calculate}
%
For instance, using methods we develop below, it turns out that the
above expression for the probability of $n$ heads in $2n$ tosses is
very close to $1/\sqrt{\pi n}$, with the approximation being more
accurate the larger $n$ is.  In fact, for $n \geq 8$, the
approximation is good to two places; for $n \geq 25$, the
approximation gives three-place accuracy.

In his book {\em Methodus differentialis\/} (1730), the British
mathematician James Stirling\index{Stirling, J.} published the
following approximation, now know as {\bf Stirling's formula}, for the
factorial operator:
$$
n!\sim \sqrt{2\pi}\,n^{n+\frac{1}{2}}\,e^{-n}.
$$

While the right-hand side may look much more complicated than the
left, think which one you would rather evaluate for, say, $n=100$. To
see how good this approximation is, here are some comparisons:
\begin{center}
\begin{tabular}{ccc}\label{table2}
&&Stirling's \\
$n$&$n!$&approximation \\ [2pt]
\hline\\[-10pt]
2&2&1.9190   \\
10 & 3,628,800   & 3,598,695.6\\
50 & $3.0414\times 10^{64}$   & $3.0363\times 10^{64}$\\
100 & $9.3326\times 10^{157}$   & $9.3248\times 10^{157}$\\
1000 & $4.02387\times 10^{2567}$   & $4.02354\times 10^{2567}$\\
10000& $2.84626\times 10^{35659}$ & $2.84624\times 10^{35659}$ \\
\end{tabular}
\end{center}

As an example of the way elementary ideas in calculus can be used to
derive powerful and subtle results, we will outline a derivation of
Stirling's approximation for $n!$.  You should write up your own
summary of this proof, filling in the gaps in the text below.
We will work in two stages.  In the first stage, we will show that
$$
n!\sim c\,n^{n+\frac{1}{2}}\,e^{-n}.
$$
for some constant $c$.  In the second stage we will show that this
constant is actually $\sqrt{2\pi}$.

\subsection{Stage One: Deriving the General Form}

We first observe that
$$
\ln (n!) = \ln 1 + \ln 2 + \ldots +\ln n.
$$
It turns out to be easier to prove things about this logarithmic
form.\index{logarithm} In fact, we will deal most easily with
$$
A_n=\ln 1 + \ln 2 + \ldots +\ln (n-1) + \dfrac{1}{2}\,\ln n.
$$
Thus $\ln(n!) = A_n+\tfrac{1}{2}\ln n \,.$ Even though $\ln 1 = 0$, it
will be useful to retain the term in the expression for $A_n$.

We will find upper and lower bounds for $A_n$ (and hence for
$\ln(n!)$\,) by approximating the area under the curve $y=\ln x$ by
certain inscribed and circumscribed trapezoids.\index{trapezoid} We
will then use these bounds to predict the asymptotic behavior of $A_n$
for large values of $n$.

\medskip
\noindent
\begin{minipage}{3.2in}
{\bf The upper bound}:\index{bound!integral} Note that if we inscribe
a trapezoid under the graph of $y=\ln x$ between $x=k-1$ and $x=k$,
its area will be $\tfrac{1}{2}(\ln (k-1) +\ln k)$. (How do we know
that the straight line connecting the points $(k-1, \ln (k-1))$ and
$(k, \ln k)$ will lie under the graph of $y=\ln x$?) The sum of the
areas of all such trapezoids from $x=1$ to $x=n$ is clearly less than
the area under the curve $y=\ln x$ over the interval $[1,n]$.
\end{minipage} \ \
\begin{minipage}{2.2in}
\special{psfile=../ps/calc2/ch121-1.ps}
\end{minipage}

\vspace{4pt}
We therefore have the inequality
$$
\dfrac{1}{2}(\ln 1 +\ln 2) +\dfrac{1}{2}(\ln 2 +\ln 3)+ \cdots + \dfrac{1}{2}(\ln(n-1)+\ln n)<\int_1^n\ln x\,dx,
$$
which is equivalent to $A_n <\dint_1^n\ln x\,dx$.

\medskip


\begin{minipage}{2.2in}
\special{psfile=../ps/calc2/ch121-2.ps}
\end{minipage}
\hfill
\begin{minipage}{3.2in}
{\bf The lower bound}: On the other hand if we draw the tangent line
to $y=\ln x$ at \mbox{$x=k$} and form the trapezoid between $x=k-.5$
and $x=k+.5$, its area will just be $\ln k$ and will be greater than
the area under the curve over the same interval. (We've used the
fact---which you should check---that the area of a trapezoid equals
the distance between the parallel sides times the distance between the
midpoints of the other two sides.)

\quad \ Adding up all such trapezoids, we get the inequality
$$
\int_{\frac{3}{2}}^n\ln x\,dx <A_n.
$$
\end{minipage} 


\vspace{6pt}

Since we know that $\int \ln x\,dx = x\ln x -x$, we can evaluate these
upper and lower bounds to conclude
$$
n\ln n-n-\dfrac{3}{2}\ln \dfrac{3}{2} + \dfrac{3}{2} <A_n<n\ln n -n +1,
$$
which in turn yields
$$
\left(n+\dfrac{1}{2}\right)\ln n-n+\dfrac{3}{2}\left(1-\ln \dfrac{3}{2}\right) <\ln n!<\left(n+\dfrac{1}{2}\right)\ln n -n +1.
$$

Pause for a moment to observe that the difference
$$
D_n = \left(n+ \dfrac{1}{2}\right)\ln n -n +1-\ln n!
$$ 
between the expressions on the two sides of the rightmost
inequality is just the accumulated error from approximating the area
under $y=\ln x$ by the inscribed trapezoids. Since the error over each
interval is always positive, $D_n$ must therefore get larger as $n$
increases, We will need this fact shortly.

Returning to our inequalities, they can finally be rewritten as
$$
\dfrac{3}{2}\left(1-\ln \dfrac{3}{2}\right) 
 < \ln n!-\left(n+ \dfrac{1}{2}\right)\ln n +n <1.
$$
Evaluating the constants, we thus have that for any value of $n$,
$$
.8918<\ln n!-\left(n+\dfrac{1}{2}\right)\ln n +n <1.
$$
If we exponentiate, this becomes
$$
2.395<\dfrac{n!}{n^{n+\frac{1}{2}}e^{-n}} <2.719,
$$
or
$$
2.395 \, n^{n+\frac{1}{2}}e^{-n} 
   < n! < 2.719 \, n^{n+\frac{1}{2}}e^{-n}.
$$

Notice that these bounds are already quite strong 
%
\mar{These estimates\\ are often adequate}	
%
and would be adequate for many estimates.  Moreover, they are true for
any value of $n$.  If we are only interested in large values of $n$,
we can do a little better.  Let
$$
\delta_n=\ln n!-\left(n+\dfrac{1}{2}\right)\ln n +n.
$$ 
Then $1-\delta_n=(n+\frac{1}{2})\ln n -n +1-\ln n!$ is just the
expression we called $D_n$ a moment ago and said had to be increasing
as $n$ increases.  But if $D_n=1-\delta_n$ is increasing, it must be
true that $\delta_n$ itself is decreasing as $n$ gets larger.  We thus
must have $1>\delta_1>\delta_2>\ldots >\delta_n \ldots >.8918$.  There
must therefore be some constant $d\geq .8918$ such that
$\lim_{n\rightarrow \infty}\delta_n=d$. Define the constant $c$ by
$c=e^d$. Then
$$
\lim_{n\rightarrow \infty}\dfrac{n!}{n^{n+\frac{1}{2}}e^{-n}} = c,
$$
which is what we mean when we write
$$
n!\sim c\,n^{n+ \frac{1}{2}}\,e^{-n}.
$$
This completes stage 1.  In stage 2 we will see that $c=\sqrt{2\pi}$.



\subsection{Stage Two: Evaluating $c$}

We will do this using an interesting result of a 17th century English
mathematician, John Wallis,\index{Wallis, J.}\index{Wallis's formula}
%
\mar{\vspace{22pt}Wallis's formula}
%
who showed that
\[
\lim_{n \rightarrow \infty} \dfrac{2}{1} \times \dfrac{2}{3}
   \times \dfrac{4}{3}\times \dfrac{4}{5}\times \dfrac{6}{5}
   \times \dfrac{6}{7}\times \cdots \times \dfrac{2n}{2n-1}
   \times \dfrac{2n}{2n+1} = \dfrac{\pi}{2}.
\]

Suppose for the moment that we had proved Wallis's formula. We can
express it in terms of factorials by noting that we can rewrite the
product of the first $n$ even numbers---$2\times 4\times 6\times
\ldots \times (2n)$---by factoring a $2$ out of each term, leaving us
\[
2\times 4\times 6\times \ldots \times (2n) = 2^n\,(n!).
\]
Similarly, we can take the product of the first $n$ odd
integers---$1\times 3\times 5\times 7\times \ldots \times
(2n-1)$---and insert the missing even terms to get
\begin{align*}
1\times 3\times 5\times 7\times \ldots \times (2n-1) 
  &= \dfrac{1\times 2\times 3\times 4\times \ldots \times (2n-1)
            \times (2n)}
      {2\times 4\times 6\times \ldots \times (2n)} \\
  &= \dfrac{(2n)!}{2^n\,(n!)}.
\end{align*}
We can thus rewrite Wallis's formula as
$$
\lim_{n \rightarrow \infty} \dfrac{\left(2^n\,n!\right)^4}{\left((2n)!\right)^2(2n+1)}= \dfrac{\pi}{2}.
$$

If we now replace all the factorials by their corresponding expressions using Stirling's approximation, we get
$$
\lim_{n \rightarrow \infty} \dfrac{2^{4n}c^4n^{4n+2}e^{-4n}}{c^2(2n)^{4n+1}e^{-4n}(2n+1)}= \dfrac{\pi}{2},
$$
which, after a great deal of cancelation, reduces to
$$\lim_{n \rightarrow \infty} \dfrac{c^2n}{2(2n+1)}= \dfrac{\pi}{2}.
$$

Now since
$$\lim_{n \rightarrow \infty} \dfrac{n}{2n+1} = \dfrac{1}{2},
$$
this reduces to
$$
\dfrac{c^2}{4} = \dfrac{\pi}{2},
$$
so
$$c^2 = 2\pi,
$$
and
$$
c=\sqrt{2\pi},
$$
as desired.

\subsubsection{Deriving Wallis's formula}
One way to derive Wallis's formula involves the integrals
	$$
	I_k = \int_0^{\pi/2}\sin^kx\,dx.
	$$
Note that $I_0>I_1>I_2>I_3>\ldots .$  Moreover, you should verify that
	$$
	I_0 = \dfrac{\pi}{2} \qquad \mbox{and} \qquad I_1=1.
	$$ 
Using the reduction formula derived in chapter~11.5 for
antiderivatives of $\sin^n{x}$, we have a similar reduction formula
for the $I_k$:
\begin{align*}
I_k &= \int_0^{\pi/2}\sin^kx\,dx  \\
  &= \left. \dfrac{-1}{k}\sin^{k-1}x\,\cos x \, \right|_0^{\pi/2}
      + \dfrac{k-1}{k} \int_0^{\pi/2} \sin^{k-2}x\,dx \\
  &=  \dfrac{k-1}{k} I_{k-2}.
\end{align*}

This in turn leads to
\[
I_k  = \begin{cases}
\dfrac{2n-1}{2n}\cdot \dfrac{2n-3}{2n-2}\cdots \dfrac{1}{2}\cdot  
    \dfrac{\pi}{2} & \mbox{if $k=2n$ is even}, \\[.15in]
\dfrac{2n}{2n+1}\cdot \dfrac{2n-2}{2n-1}\cdots \dfrac{2}{3}\cdot  1
  & \mbox {if $k=2n+1$ is odd}.
\end{cases}
\]
Further, note that
	$$
	I_{2n+2}/I_{2n}= \dfrac{2n+1}{2n+2},
	$$
which has the limit 1 for large $n$.  Since
$I_{2n}>I_{2n+1}>I_{2n+2}$, it follows that $I_{2n+1}/I_{2n}$
approaches 1 for large $n$.  But this gives us
\begin{align*}
1 &= {\displaystyle\lim_{n \rightarrow \infty}I_{2n+1}/I_{2n}} \\[4pt]
  &= {\displaystyle  \lim_{n \rightarrow \infty}
    	\left(\dfrac{2n}{2n+1}\cdot \dfrac{2n-2}{2n-1}\cdots 
		\dfrac{2}{3}\cdot  1\right)
	\div 
\left(\dfrac{2n-1}{2n}\cdot \dfrac{2n-3}{2n-2}\cdots \dfrac{1}{2}\cdot 
	\dfrac{\pi}{2}\right)} \\[4pt]
   &= {\displaystyle\lim_{n \rightarrow \infty}
	\dfrac{2n}{2n+1}\cdot \dfrac{2n}{2n-1}\cdot \dfrac{2n-2}{2n-1}\cdots 
	\dfrac{2n-2}{2n-3} \cdots \dfrac{2}{3}\cdot \dfrac{2}{1}\cdot 
	\dfrac{2}{\pi}}
\end{align*}
If we multiply both sides of this equation by $\pi/2$, we get Wallis's formula.

\subsubsection{Further refinements}

Using even more careful methods of analysis, 
%
\mar{Some refinements}
%
it is possible to improve on Stirling's approximation and derive
approximations like
$$
n!\sim \sqrt{2\pi}\,n^{n+\frac{1}{2}}\,
     e^{-n+\frac{1}{12n}-\frac{1}{360n^3}+\frac{1}{1260n^5}-\cdots}.
$$ 
If we use this expression to approximate $1000!$, for instance, our
result is accurate for the first 24 digits.

While this approximation and Stirling's original one are good in the
sense that they give more and more accurate digits the larger $n$
gets---so that the {\em ratio} of $n!$ to either approximation goes to
1 as $n$ gets large---they are bad in the sense that the {\em
  difference} between $n!$ and either approximation becomes infinite
as $n$ gets large.

\subsection{The Binomial Distribution}\index{binomial probability distribution}\index{probability distribution!binomial}

One of the most frequently encountered concepts in probability theory
is the {\bf binomial probability distribution}. Suppose we repeat a
certain experiment---flipping a penny, rolling a single die, mating a
pair of fruit flies, feeding cholesterol to a lab rat---over and over.
Suppose further that there is some outcome we are looking
for---getting heads, rolling a 2, getting a red-eyed offspring,
developing liver cancer in the rat---in each experiment.  If $p$ is
the probability $p$ of obtaining the looked-for outcome in any one
experiment, denote by $P(n, k, p)$ the probability of the outcome
happening exactly $k$ times in $n$ experiments.  It turns out that
$$
P(k,n, p)=\dfrac{n!}{k!(n-k)!}p^k(1-p)^{n-k}.
$$

\medskip
\noindent
{\bf Example 1}\quad How likely is it to get four 2's if we roll
twelve dice?  The probability of getting a 2 by throwing one die is
$\dfrac{1}{6}$.  Therefore the answer to the question is
$$
P(12, 4, \tfrac{1}{6})= \dfrac{12!}{4!\, 8!}
   \left(\dfrac{1}{6}\right)^4\left(\dfrac{5}{6}\right)^8=.0888281
$$ 
---we should get exactly four 2's slightly less frequently than once
out of every 11 times we roll twelve dice.

\medskip
\noindent
{\bf Example 2}\quad What is the probability of getting exactly 47
heads if we flip 100 pennies? Since the probability of getting heads
on a single toss of a penny is $\frac{1}{2}$,
$$
P(100, 47, \tfrac{1}{2}) 
 = \dfrac{100!}{47!\,53!}\left(\dfrac{1}{2}\right)^{100}=.0665905,
$$
---on the average, if we flip 100 pennies, we should get 47 heads
about once out of every 15 times.

\medskip

The second example demonstrates the fact that calculating binomial
probabilities can get very messy very quickly.  Several of the
exercises are designed to show how Stirling's formula can give us
quick estimates that are easy to calculate and work with.

\subsection{Exercises}

\vspace{-10pt}

\num Go through the derivation in this section and find several
passages that seem to you to go a bit fast or skip over details.
Rewrite these sections to make them clearer and more complete.

\num Confirm the values given in the table on page~\pageref{table2}
for the approximations of $100!$ and $1000!$ that Stirling's formula
produces.

\num {\bf Rate of growth of $n!$}\quad Factorials get very large very
rapidly.  The purpose of this exercise is to develop a sense of just
how rapidly $n!$ grows by comparing it to exponential functions.

Let $N$ be some integer $>1$, and consider the sequence $a_1$, $a_2$,
$a_3$, \ldots defined by
$$
a_n=\dfrac{N^n}{n!}.
$$

\sub Show that $a_k = \dfrac{N}{k}a_{k-1}$, and conclude that 
\begin{center}
\begin{tabular}{c}
if $k<N$\quad then\quad $a_{k-1}<a_k$; \\
if $k>N$\quad then\quad $a_{k-1}>a_k$; \\
if $k=N$\quad then\quad $a_{k-1}=a_k$. \\
\end{tabular}
\end{center}
We thus have a sequence that increases for a while: 
$$
a_1<a_2<\cdots a_{N-1}=a_N,
$$ 
and then decreases forever after:
$$
a_N>a_{N+1}>a_{N+2}>\cdots .
$$

\sub If $k>2N$ show that $a_k<.5a_{k-1}$.  Hence conclude that
${\displaystyle \lim_{n\rightarrow\infty}a_n=0}$.

\sub Use Stirling's approximation to show that
$$
a_N\approx \dfrac{e^N}{\sqrt{2\pi N}}.
$$
Calculate the values of this expression for $N=10$ and $N=100$ to get
an idea of how large the sequence $\{a_n\}$ can get.  This shows that
{\em for a while}, the exponential series $\{N^n\}$ can get large much
more rapidly than the series $\{n!\}$.

\sub Show that $a_n<1$ if $n>eN$.  This gives an upper bound on how
long it takes the factorials to catch up with the exponentials.

\num If $n\geq 5$, then $n!$ terminates in a certain number of zeroes.
For instance, $5!=120$ ends in one zero, $23!=25852016738884976640000$
ends in four zeroes, and so on.  How many zeroes are there at the end
of $1000!$\,?

\subsubsection{The binomial distribution}

\vspace{-10pt}

\num The formula for the binomial distribution gives us that the
probability of getting exactly $n$ heads in $2n$ flips of a coin is
	$$
	\dfrac{(2n)!}{(n!)^2}\left(\dfrac{1}{2}\right)^{2n}.
	$$  
Show using Stirling's formula this can be approximated by
$$
\dfrac{1}{\sqrt{\pi n}}.
$$ 
Use this approximation to find the probability of getting 50 heads out
of 100 tosses of a coin.  If you have a computer or calculator which
can compute factorials, use the original binomial distribution formula
to calculate the exact probability of getting 50 heads and compare the
answers.

\num More generally, if we try a certain experiment $n$ times with a
probability $p$ of success each time, the most likely number of
successes is $k=np$.  (Assume that $p$ is a fraction and $n$ is such
that $n\cdot p$ is an integer.)  Use Stirling's approximation to show
that the probability of getting exactly $np$ successes is
$$
P(n, np, p)\approx \dfrac{1}{\sqrt{2\pi np(1-p)}}.
$$
Is this consistent with the answer to the previous exercise?

\num {\bf One-dimensional random walk}\quad An important class of
problems, including {\bf diffusion} and {\bf Brownian motion} involve
the long-term behavior of particles moving randomly.  We will look at
the simplest case of such problems.  A particle starts at the origin
on a line and at each stage moves one unit to the right or one unit to
the left, being equally likely to do either.  What can we say about
where the particle will be after $n$ steps?  In this problem we will
use Stirling's formula to develop some useful insights into this
question.

\sub Explain why the particle will be $r$ units to the right of the
origin after $n$ steps if and only if it has moved to the right
$k=(n+r)/2$ times and to the left $n-k=(n-r)/2$ times. Explain why it
could never be 3 units or 7 units to the right after 100 steps.

\sub Using the same symbols as in part (a), show that the probability
of the particle's being exactly $r$ units to the right after $n$ steps
is
$$
\dfrac{n!}{k!(n-k)!}\left(\dfrac{1}{2}\right)^n.
$$

\sub Use Stirling's formula to show that this probability of being $r$
units to the right after $n$ steps is approximately
$$
\dfrac{\sqrt{2}}{\sqrt{\pi n} (1+(r/n))^{\frac{1}{2}(n+r+1)} \, 
     (1-(r/n))^{\frac{1}{2}(n-r+1)}}.
$$

\sub To simplify the denominator of this fraction, recall the Taylor
series approximation for $\ln(1+x)$:
$$
\ln(1+x)=x-\dfrac{x^2}{2} +\cdots .
$$
Hence, if $r$ is much smaller than $n$, $\ln(1+r/n)$ can be
approximated by $r/n-r^2/(2n^2)$, and $\ln(1-r/n)$ can be approximated
by $-r/n-r^2/(2n^2)$. By ignoring all powers of $r$ greater than the
second, conclude that
$$
(1+(r/n))^{\frac{1}{2}(n+r+1)}\,(1-(r/n))^{\frac{1}{2}(n-r+1)}
    \approx e^{r^2/(2n)},
$$
so that the probability of being $r$ units to the right after $n$ steps is
$$
\sqrt{\dfrac{2}{\pi n}} \, e^{-r^2/2n}.
$$

\sub Explain how we can get the answer to exercise~5 as a special case
of the result just obtained in part (d).

\sub Using the approximation from part (d), calculate the probability
that after 100 steps the particle will be {\em no more} than 5 units
away from the starting point to either the right or the left.
Remember that after 100 steps it is impossible to be an odd number of
units away from the starting point.  The exact probability is
\[
\sum_{k=48}^{52} \dfrac{100!}{k!(100-k)!}
  \left(\dfrac{1}{2}\right)^{100} = .382701.
\]


\newpage
%%%%%%%%%%%%%%%%%%%%%%%%%%%%%%%%%%%%%%%%%%%%%%%%%%%%%%%%%%%%%%%%%%%%%%%
%                                                                     %
%                              Section 2                              %
%                                                                     %
%%%%%%%%%%%%%%%%%%%%%%%%%%%%%%%%%%%%%%%%%%%%%%%%%%%%%%%%%%%%%%%%%%%%%%%


\section{The Poisson Distribution}
\index{Poisson distribution}\index{probablity distribution!Poissson}

\subsection{A Linear Model for $\alpha$-Ray Emission}
\index{$\alpha$-ray}
\index{alpha ray}

When a radioactive element decays, we know from the study of
differential equations in chapter 4 that the amount $A(t)$ of
radioactive material present at time $t$ satisfies the differential
equation
$$
A' = -kA,
$$
where $k > 0$ is the decay constant.  If $A_0$ is the amount present
at time $t = 0$, then the solution is
$$
A(t) = A_0e^{-kt}.
$$

The time $T$ it takes for a given amount of radioactive material to
decay to half the starting quantity is known as the {\bf half
  life}\index{half life} of the element.  Since, by definition,
$A(T)=.5\,A(0)=.5\,A_0$, we must have
$$
e^{-kT}=\dfrac{1}{2},
$$
which leads to
$$
kT=\ln{2}
$$
and 
%
\mar{The relation between the half-life and\\ the decay constant}
%
therefore
$$
T=\dfrac{\ln{2}}{k}.
$$

Suppose, for example, that we have a sample of
polonium,\index{polonium} which is a radioactive isotope of
radium.\index{radium} The decay constant of polonium is $k = .500865$
\% per day, and thus its half life is
$$
T = \dfrac{{\ln2}}{k} = \dfrac{{\ln2}}{.00500865} 
   = 138.39\ \hbox{days}.
$$
By linear approximation, $A(t)$ is closely approximated by a linear
function for short intervals of time.  Because polonium has a
half-life of 138.39 days, a ``short time'' means several hours in this
case.  Thus, if we spend an afternoon in a laboratory studying the
decay of polonium, we can assume that $A(t)$ is linear.

When polonium decays, it produces various sorts of
radiation,\index{radiation} including $\alpha$-rays (``alpha rays'').
Using a scintillation counter, one can determine the number of rays
emitted in given directions:

\vspace*{2in}
%\special{psfile=ch122-1a.ps}
%\special{psfile=ch122-1b.ps}

\noindent 
A setup like this will count a fixed percentage of the total number of
$\alpha$-rays emitted.  Since our model of decay\index{radioactive
  decay model} is linear, it follows that the number of $\alpha$-rays
detected should be a linear function of time.  If we start counting at
time $t = 0$, the number of particles observed will have a
straight-line graph:\label{linear}

\vspace*{2in}
%\special{psfile=ch122-2.ps}

\enlargethispage*{10pt}
In the early 20th century, researchers like Marie Curie\index{Curie,
  M.} and Ernest Rutherford\index{Rutherford, E.} did numerous studies
of the $\alpha$-rays emitted by polonium.  For example, in 1911,
Rutherford, Geiger\index{Geiger} and Bateman\index{Bateman} counted
the number of $\alpha$-rays detected in a 7.5-second time period.
They repeated their experiment 2608 times and detected a total of
10,097 $\alpha$-rays.  This is an average of
$$
\dfrac{10097}{2608} = 3.8715\ \hbox{$\alpha$-rays per 7.5-second period},
$$
so the number of $\alpha$-rays per second is 
$$
\dfrac{3.8715}{7.5} = .5162\ \hbox{$\alpha$-rays per second}.
$$
Thus the straight line in the above graph has slope .5162.

\newpage

This model of $\alpha$-ray production has several problems. First, it
predicts the existence of fractional $\alpha$-rays, which makes no
sense---the number detected is always a nonnegative integer. To remedy
this, we can modify our model as follows:

\vspace*{2in}
%\special{psfile=ch122-3.ps}


\noindent Notice that the graph is now a step function\index{function!step}.  It shows that we should see a new $\alpha$-ray every $1/.5162 = 1.937$ seconds. This model also has the following consequence: if we observe the number of $\alpha$-rays produced in a 7.5-second interval, then we will always see 3 or 4 particles:
\vspace*{2.3in}
%\special{psfile=ch122-4.ps}


\noindent 
As the picture indicates, whether we get 3 or 4 depends on where the
interval starts.  Now comes the serious problem: this prediction is
{\em inconsistent\/} with the experimental data collected by
Rutherford\index{Rutherford, E.} and the others in 1911.  For example,
in 57 of the 2608 times they ran the experiment, {\em no\/}
$\alpha$-rays were observed, while in 139 cases, 7 $\alpha$-rays were
observed.  Here are the complete data of the experiment:
\begin{center}
	\begin{tabular}{cc}\label{data}
	number of & number of \\
	$\alpha$-rays observed & occurrences \\
	$n$ & $N_n$ \\ [2pt] \hline
		\begin{tabular}{r}
		\rule{0pt}{14pt} 0 \\
		1 \\
		2 \\
		3 \\
		4 \\
		5 \\
		6 \\
		7 \\
		8 \\
		9 \\
		10 \\
		11 \\
		12 \\
		13 \\
		14 
		\end{tabular}
		&
		\begin{tabular}{r}
		\rule{0pt}{14pt} 57\\
		203\\
		383\\
		525\\
		532\\
		408\\
		273\\
		139\\
		 45\\
		 27\\
		 10\\
		  4\\
		  0\\
		  1\\
		  1
		\end{tabular}
	\\ [6pt] \hline
	Total & \rule{0pt}{14pt} 2608 \rule{6pt}{0pt}
	\end{tabular}
\end{center}

It follows that the linear model of $\alpha$-ray emission doesn't
apply to time intervals of length 7.5 seconds.  This is a common
occurrence---a model may work nicely over a certain range, but outside
of that range, its answers may be meaningless.  The problem in our
case comes from the random nature of radioactive decay.  In fact,
there are two sources of randomness to deal with: the {\em time\/}
when a polonium atom decays is random, and the {\em direction\/} in
which it then emits an $\alpha$-ray is also random (this affects us
since the scintillation counter only detects emissions in certain
directions).  We need to modify our model to take the randomness into
account, and this is where probability enters in.

\subsection{Probability Models}
\index{model!probability}


The basic idea of probability theory is that the outcome
%
\mar{Randomness has structure} 
%
of a certain event can be unpredictable in the individual instance but
predictable on the average.  Throwing dice and tossing a coin are
familiar examples.  In this section, we will show how the Poisson
probability distribution gives an excellent model of the $\alpha$-ray
experiment described above.

\subsubsection{The definition of probability}
\index{probability}

We will let $p_n$ denote the probability of observing exactly $n$
$\alpha$-rays in a 7.5-second time interval.  By this statement, we
mean the following.  Suppose we run the experiment $N$ times, where
$N$ is large.  Let $N_n$ be the number of times we observed $n$
$\alpha$-rays.  Then the ratio $N_n/N$ is the frequency with which
this outcome occurs.  Now imagine $N$ getting larger and larger. Being
``predictable on the average'' means that the ratios $N_n/N$ approach
a fixed number, that is, the limit $\lim_{N\to\infty}N_n/N$ exists.
We then {\em define\/} this number to be the probability $p_n$.  Thus
\[ 
p_n = \lim_{N\to\infty}\dfrac{N_n}{N}.
\]
For example, the data presented on page~\pageref{data} were obtained
from $N = 2608$ repetitions of our experiment.  From the table given
there, we see that 0 $\alpha$-rays were observed 57 times.  This means
$N_0 = 57$, and thus the probability of detecting 0 $\alpha$-rays is
\[ 
p_0 \approx \dfrac{N_0}{N} = \dfrac{57}{2608} = .0218.
\]
Similarly, we can approximate $p_1$, $p_2$, etc.,~using the data in
the table.  Our goal is to describe these probabilities $p_0$, $p_1$,
\ldots \,..  Ideally, we would like to have a way of determining the
numbers $p_0, p_1, p_2, \ldots$ ``before the fact."

\subsubsection{Some properties of probabilities}

In any introductory course on probability, one learns certain basic
principles for working with probabilities.  We will give examples to
illustrate some of these principles, and more examples may be found in
the exercises.

For our purposes, we will be working in the following setting.  There
is a certain {\bf experiment} being performed.  This might consist of
flipping
%
\mar{The general context}
%
a coin and noting which side comes up, or running a survey asking
people at random their opinions about a certain TV show, or, in our
case, counting the number of $\alpha$-rays detected in a 7.5-second
interval.  Moreover, there is a {\bf discrete} set of possible
outcomes of the experiment.  That is, the possible outcomes can be
listed in a sequence $O_1$, $O_2$, $O_3$, \ldots\,.  In some cases,
like throwing a pair of dice, this list might be finite.  In other
cases, like our $\alpha$-ray experiment, the list might be infinite.
What is ruled out are experiments like choosing a person at random and
measuring the person's height---there is a continuum of possible
outcomes here which cannot be listed in the way we've specified.
Moreover, there should be a {\bf probability} assigned to each
outcome, with the outcome $O_n$ having probability $p_n$..  Finally,
the possible outcomes should be {\bf disjoint}---two different
outcomes can't both result from a single experiment.  Thus if we are
examining the attributes of a group of people, ``being male'' and
``having green eyes'' would not be acceptable outcomes in our sense
unless we somehow knew in advance that there were no green-eyed males
in the group.

Knowing the probabilities $p_0,p_1,\ldots$ of the possible outcomes
allows us to compute other, possibly more complicated probabilities.
This brings in the concept of an {\bf event}, which is basic to
probability.  In the case of our $\alpha$-ray experiment, here are
some examples of events:
\begin{itemize}

\item Detecting 3 $\alpha$-rays.

\item Detecting 2 or 4 $\alpha$-rays.

\item Detecting an odd number of $\alpha$-rays.
\end{itemize}
In general, an {\bf event}\index{event!in probability} is a
subcollection of the possible outcomes.  
\mar{\vspace*{18pt} The addition rule\\ for probabilities}
\begin{quote}  
{\bf Rule 1\quad The {\bf probability of an event}\index{probability!of an event} is simply the sum of the probabilities of its component outcomes.}\label{rule1}
\end{quote}
Thus, for the events just described, we have:
\begin{itemize}

\item The probability of detecting 3 $\alpha$-rays is $p_3$;

\item The probability of detecting 2 or 4 $\alpha$-rays is $p_2+p_4$;

\item The probability of detecting an odd number of $\alpha$-rays is
  the infinite sum
\[
p_1+p_3+p_5+p_7+\cdots 
\]
(since an odd number of $\alpha$-rays means that 1 or 3 or 5 or 7
etc.\ have been detected).
\end{itemize}

Another important property of probabilities follows directly from Rule 1:
\begin{quote}
{\bf Rule 2\quad The sum of the probabilities of all possible outcomes is~1:
\[
 \sum_{k=0}^\infty p_k=1.
\]
}
\end{quote}
The reason for this is that the list of outcomes was stipulated to be
the list of {\em all possible} outcomes. Hence the event consisting of
all these outcomes is bound to occur every time---its probability is
1.

A third rule we will need relates the probabilities of {\bf
  independent} events.  Two events are independent if the occurrence
or non- occurrence of one of the events has no impact on the
probability of the second event occurring.  For instance, suppose we
are examining a group of people.  Consider the following events which
may or may not occur each time we look at a person:
\begin{enumerate}

\item The person is female;

\item The person has green eyes;

\item The person is over 5'7" tall.
\end{enumerate}
We would expect the first and second events to be independent, and
also the second and third, but not the first and third.
\mar{\vspace*{18pt} The product rule\\ for probabilities}
\begin{quote}
{\bf Rule 3\quad The probability that two or more independent events
  all occur is the product of their separate
  probabilities.}\label{rule3}
\end{quote}
Thus, for example, suppose that in our hypothetical group of people
$\frac{1}{2}$ are female, $\frac{1}{8}$ are green-eyed, and
$\frac{1}{3}$ are taller than 5'7".  We might then expect roughly
$\frac{1}{24}$ of them to be green-eyed {\em and} over 5'7", but we
would have no particular reason to expect that $\frac{1}{6}$ of them
are females taller than 5'7".

A final rule that is often useful is
\mar{\vspace*{16pt} The probability\\ that something\\ doesn't happen}
\begin{quote}
{\bf Rule 4\quad If a certain event has a probability \boldmath $p$ of
  happening, then the probability that the event doesn't take place is
  $1-p$.}\label{rule4}
\end{quote}
For example, in our group of people, we would expect $\frac{2}{3}$ of
them to be less than 5'7" tall, $\frac{7}{8}$ of them to have eyes
colored something other than green, etc.

\subsubsection{The notion of a probability model}

A {\bf model}\index{model} is a mathematical picture of a real-life
phenomenon.  We have seen that dynamical systems can be used to create
models of physical situations.  Another type of mathematical model is
a {\em probability model}.  In general, a {\bf probability
  model}\index{probability model} for an experiment with a finite
number of outcomes is {\em a listing of all possible outcomes and an
  assignment of probabilities to each outcome so that their sum is 1}.
In order that the probability model be a good picture of reality, we
ask that the {\em probability assigned to an outcome should be the
  relative frequency with which that outcome would appear if the
  experiment were duplicated independently a large number of times}.

As an example, a probability model for one toss of a fair die consists
of a list of all possible outcomes, namely 1, 2, 3, 4, 5, 6, and an
assignment of a probability to each, namely $\frac{1}{6}$,
$\frac{1}{6}$, $\frac{1}{6}$, $\frac{1}{6}$, $\frac{1}{6}$,
$\frac{1}{6}$, respectively.  We assign the number $\frac{1}{6}$ to
each outcome because we expect that if the the experiment were
repeated (that is, if the die were tossed) a large number of times,
then any particular outcome (3, say) would occur about one sixth of
the time.  Another probability model for the experiment consisting of
a toss of a die might be a list of all outcomes, again 1, 2, 3, 4, 5,
6, together with an assignment of the numbers $\frac{1}{2}$, 0,
$\frac{1}{6}$, 0, 0, $\frac{1}{3}$ to 1, 2, 3, 4, 5, 6, respectively.
This is a probability model, because the numbers we have assigned add
to 1, but it certainly does not model very well the throw of a fair
die.

We would like to set up a probability model for our experiment with
$\alpha$-rays.  The outcomes are 0, 1, 2, 3, 4, ... where, for
example, the number 5 labels the outcome in which we observe 5
$\alpha$-rays in our 7.5-second interval.  The total number of
outcomes is equal to the number of $\alpha$-rays that we could
conceivably see in a 7.5-second interval.  Since it is conceivable
(but extremely unlikely) that every atom in the sample could decay and
emit an $\alpha$-ray in the direction of the scintillation counter in
one 7.5-second interval, we could conceivably see as many
$\alpha$-rays as there are atoms in the sample. This number is so
large that we can think of it as infinite.  To have a probability
model, we need to assign numbers $p_0$, $p_1$, $p_2$, \ldots to the
outcomes 0, 1, 2, \ldots, respectively, so that $p_0 + p_1 + p_2 +
\cdots = 1$.  For the model to be reasonable, we would like each $p_n$
to be approximately equal to the corresponding number $N_n/N$ observed
by Rutherford, Geiger, and Bateman.

\newpage

\subsection{The Poisson Probability Distribution}


\subsubsection{The Poisson model of $\alpha$-ray emission}
\index{Poisson model}

To describe the probabilities $p_0$, $p_1$, \ldots, $p_n$, \ldots that
we will observe 0, 1, \ldots, $n$, \ldots $\alpha$-rays in a
7.5-second interval for our $\alpha$-ray experiment, we use the {\bf
  Poisson probability distribution}\index{Poisson distribution}
\[
p_n = \dfrac{{\lambda^n e^{-\lambda}}}{{n!}},
\] 
where $\lambda$ is a number yet to be determined, and $n!$ is the
familiar \textbf{\boldmath $n$-factorial}
function,\index{factorial}\index{$n"!$}
\[
n! = \begin{cases}
n \cdot (n-1) \cdot (n-2) \cdot \cdots \cdot 3 \cdot 2 \cdot 1 
    & \text{if $n > 0$}, \\ 
1 & \text{if $n = 0$}.
\end{cases}
\]
Thus the first few Poisson probabilities are:
\[
p_0 = e^{-\lambda}, \qquad
p_1 = \lambda e^{-\lambda}, \qquad
p_2 = \dfrac{\lambda^2e^{-\lambda}}{2}, \qquad
p_3 = \dfrac{\lambda^3e^{-\lambda}}{6}.
\]

Note that this assignment does indeed give us a probability model, because
\begin{align*}
p_0 + p_1 + p_2  +p_3 + \cdots 
  &=  \dfrac{\lambda^0}{0!}e^{-\lambda} + \dfrac{\lambda}{1!} e^{-\lambda} 
      + \dfrac{\lambda ^2}{2!}e^{-\lambda} 
      + \dfrac{\lambda ^3}{3!}e^{-\lambda} +\cdots \\
  &= e^{-\lambda}\left( 1 + \dfrac{\lambda}{1!} + \dfrac{\lambda ^2}{2!} 
       + \dfrac{\lambda ^3}{3!} + \cdots \right) \\
  &= e^{-\lambda}\cdot e^{\lambda}\\
  &=  1. 
\end{align*}
(The transition from the second line to the third uses the fact that
the expression in parentheses is just the Taylor series for
$e^{\lambda}$.)

We will shortly derive the Poisson distribution from basic principles.
For the moment, though, we will assume that the probabilities $p_0$,
$p_1$, $\ldots$ for $\alpha$-ray emission are given by the above
formulas, where we still need to choose an appropriate value for the
parameter $\lambda$.  The key to determining $\lambda$ is the notion
of {\bf expectation},\index{expectation} which for us will mean the
average number of $\alpha$-rays observed in a 7.5-second interval.

Suppose we repeat our experiment $N$ times.  As usual, we let $N_n$
denote the number of times exactly $n$ $\alpha$-rays were observed.
Then the total number of $\alpha$-rays observed in the $N$ experiments
is
\[
0\cdot N_0 + 1\cdot N_1 + 2\cdot N_2 + 3\cdot N_3 + \cdots\ .
\]
Then the ``average number of $\alpha$-rays observed in a 7.5-second
interval'' means the limit
\[
E = \lim_{N\to\infty} 
    \dfrac{0\cdot N_0 + 1\cdot N_1 + 2\cdot N_2 + 3\cdot N_3 + \cdots}
    {N}.
\]
This limit is called the {\bf expected value}\index{expected value} or
{\bf expectation} (which explains why it is denoted $E$).

We claim that for the Poisson distribution, the expected value $E$ is
exactly the number~$\lambda$.  To see this, notice that the above
limit can be written in the form
\[
E = \lim_{N\to\infty} \left( 0 \cdot \dfrac{N_0}{N} + 1\cdot \dfrac{N_1}{N}
      +  2\cdot \dfrac{N_2}{N} + 3\cdot \dfrac{N_3}{N} + \cdots\ \right).
\]
Since we defined
\[p_n = \lim_{N\to\infty} \dfrac{N_n}{N}\ ,\]
it follows 
%
\mar{\vspace{12pt}The general formula for the expected value in a probability model}
%
that we get the following formula for the expectation:
\[
E = 0\cdot p_0 + 1\cdot p_1 + 2\cdot p_2 + 3\cdot p_3 + \cdots =
        \sum_{n=0}^\infty np_n.
\]
(Note that this equality is true for {\em any} probability model, not
just the one we are considering)

Substituting in the values of $p_n$ given by the Poisson distribution,  we have\index{summation}
\begin{align*}
E &= 0\cdot e^{-\lambda} + 1\cdot \lambda e^{-\lambda} + 
        2\cdot \dfrac{\lambda ^2}{2!} e^{-\lambda} + 
        3\cdot \dfrac{\lambda ^3}{3!}e^{-\lambda} + \cdots \\
  &= \sum_{n=0}^\infty n \dfrac{\lambda ^n}{n!}e^{-\lambda}\\
  &= \lambda e^{-\lambda}\sum_{n=1}^\infty \dfrac{\lambda ^{n-1}}{(n-1)!},
\end{align*}
where we pulled the common factor $\lambda e^{-\lambda}$ outside the
summation, noted that the term in the summation corresponding to $n=0$
is 0, and observed that
$$
\dfrac{n}{n!} = \dfrac{n}{n(n-1)\cdot\cdots\cdot2\cdot 1} 
   = \dfrac{1}{(n-1)!}.
$$ Letting $k=n-1$, we have (again recognizing the Taylor series for
   $e^{\lambda}$)
\[
E =  \lambda e^{-\lambda}\sum_{k=0}^\infty \dfrac{\lambda ^k}{k!}
  = \lambda e^{-\lambda}e^{\lambda}
  = \lambda.
\]
This proves that the expected value is $\lambda$ as claimed.

Now that we know how to interpret $\lambda$, it is easy to determine what it should be for the $\alpha$-ray experiment.  The data given on page~\pageref{data} covered $N = 2608$ repetitions of the experiment, with
$$
0\cdot N_0 + 1\cdot N_1 + \cdots = 10097
$$
in this case. Thus
\[
\dfrac{0\cdot N_0 + 1\cdot N_1 + \cdots}{N} = \dfrac{10097}{2608}
        = 3.8715
\]
is an approximation of the expected value $\lambda$. However, since
this is the only information about $\lambda$ we have, we will let
$\lambda = 3.8715$. Using this value of $\lambda$, we can then compare
the frequencies predicted by the Poisson distribution to the actual
data from on page~\pageref{data}:
\begin{center}\begin{small}
\begin{tabular}{ccccc}
	number of & number of & probability  & 
		Poisson & Poisson \\
	$\alpha$-rays observed & occurrences & 
		approximation & probability & prediction \\
	$n$ & $N_n$ & $N_n/N$ & $p_n$ & $2608\,p_n$ \\[2pt]
		\hline
	\begin{tabular}{r}
	\rule{0pt}{14pt} 0 \\
		1 \\
		2 \\
		3 \\
		4 \\
		5 \\
		6 \\
		7 \\
		8 \\
		9 \\
		10 \\
		11 \\
		12 \\
		13 \\
		14 
	\end{tabular}
	&
	\begin{tabular}{r}
	\rule{0pt}{14pt} 57  \\
		203 \\
		383 \\
		525 \\
		532 \\
		408 \\
		273 \\
		139 \\
		 45 \\
		 27 \\
		 10 \\
		  4 \\
		  0 \\
		  1 \\
		  1 
	\end{tabular}
	&
	\begin{tabular}{l}
		.021855 \rule{0pt}{14pt} \\
		.077837 \\
		.146855 \\
		.201303 \\
		.203398 \\
		.156441 \\
		.104677 \\
		.053297 \\
		.017254 \\
		.010352 \\
		.003834 \\
		.001533 \\
		.000000 \\
		.000383 \\
		.000383 
	\end{tabular}
	&
	\begin{tabular}{l}
		.020827 \rule{0pt}{14pt} \\
		.080632 \\
		.156083 \\
		.201426 \\
		.194955 \\
		.150953 \\
		.097402 \\
		.053870 \\
		.026070 \\
		.011214 \\
		.004341 \\
		.001528 \\
		.000492 \\
		.000146 \\
		.000040 
	\end{tabular}
	&
	\begin{tabular}{r}
		\rule{0pt}{14pt} 54.3\\
		210.3\\
		407.1\\
		525.3\\
		508.4\\
		393.7\\
		254.0\\
		140.5\\
		 68.0\\
		 29.2\\
		 11.3\\
		  4.0\\
		  1.3\\
	  	   .4\\
		   .1
	\end{tabular}
	\\ [6pt] \hline
	Totals & 2608 & 1 & 1 & \rule{0pt}{14pt} 2608
\end{tabular}\end{small}
\end{center}
The Poisson model agrees nicely\index{prediction} with the data since
for each $n$, $N_n/N$ and $p_n$ are reasonably close.  Notice that we
shouldn't expect perfect agreement since $N_n/N$ is only an
approximation to $p_n$.  We would expect these approximations to get
better as we take larger values of $N$.

Look at the last column, labelled ``Poisson prediction.''  The numbers
here are the Poisson probabilities multiplied by $N = 2608$, and they
represent the ``ideal'' number of occurrences.  This makes it easier
to compare the model to the data.  For example, the graph below plots
the number of occurrences, both actual and predicted.  The circles are
the experimental data, while the line-segment graph connects the
Poisson predictions.

Although the model seems to fit the data nicely, we should point out
that there are statistical tests which can be used measure the fit
more precisely.  These tests are part of the material covered in
courses in probability and statistics.

\vspace*{3in}
%\special{psfile=ch122-5.ps}

A final and very important point to make concerns the number of
$\alpha$-rays observed over a long period of time.  Our particular
Poisson model with $\lambda=3.8715$ only works for a 7.5-second
interval.  What happens if we count $\alpha$-rays over a longer time
period?  For simplicity, assume that we have a time interval of length
$T$ which is a multiple of 7.5 seconds, so that $T = 7.5\,N$ for some
large integer $N$.  We can regard this as running our 7.5-second
experiment $N$ consecutive times. Thus the ratio
\[
\dfrac{\hbox{total number of $\alpha$-rays observed}}{N}
\]
is an approximation to the expected value $\lambda = 3.8715$.  It
follows that
\[
\hbox{total number of $\alpha$-rays observed} \approx 3.8715\,N
   = \dfrac{3.8715}{7.5}7.5\,N
   = .5162\,T.
\]
This shows that for large time intervals, we recover the linear model
of $\alpha$-ray emissions discussed on page~\pageref{linear}.  Thus
our probabilistic model is consistent with what we did earlier and yet
allows us to describe what happens when the linear model breaks down.

\subsubsection{Derivation of the Poisson model}
\index{Poisson model!derivation}

In the previous discussion we simply assumed that the $\alpha$-ray
probabilities were given by the Poisson distribution, and found that
the Poisson probabilities agreed with the experimental data. Let's see
where the Poisson formulas come from.  It turns out that we can derive
the Poisson probabilities $p_n$ from the following assumptions:
\begin{itemize}

\item We have an extremely large number $M$ of polonium atoms;

\item Each atom has a small but equal probability of emitting an
  $\alpha$-ray that is detected by our scintillation counter in a
  7.5-second period;

\item Observing an $\alpha$-ray from a given atom is independent of
  observing an $\alpha$-ray from any other atom.
\end{itemize}

Now suppose that we see an average of $\lambda = 3.8715$ $\alpha$-rays
in a 7.5-second period.  Because the number of atoms $M$ is large (in
the Rutherford-Geiger-Bateman experiment $M > 10^{18}$), then the
probability that a single fixed atom emits an $\alpha$-ray detected by
our scintillation counter in a given time period is very close to
$\lambda/M$.  The probability that the single atom does not emit a
detected $\alpha$-ray in the period is then $1-\lambda/M$ (by Rule 4,
page~\pageref{rule4}). Thus, the probability $p_0$ that none of the
$M$ atoms emits an $\alpha$-ray in the 7.5-second period is
$(1-\lambda/M)^M$ (by Rule 3, page~\pageref{rule3}).

The fact that $M$ is so large allows us to make a simplifying
approximation.  Recall that for any value of $x$, positive or
negative,
$$
e^x = \lim_{n\rightarrow\infty}\left(1+\dfrac{x}{n}\right)^n.
$$
Therefore  
$$
p_0 = \left(1-\dfrac{\lambda}{M}\right)^M \approx 
	e^{-\lambda}.
$$ 
You might calculate some sample values for various values of $M$ to
see how good this approximation is.

To derive the values of $p_k$ for $k>1$, we need a slight improvement
on the estimate we just made.  Let $k$ be a relatively small (compared
to the size of $M$) number. Then
$$
\left( 1-\dfrac{\lambda}{M}\right)^{M-k}
  = \left(\left( 1- \dfrac{\lambda}{M}\right)^M\right)^{\frac{M-k}{M}}\ .
$$
Now if $M$ is very large compared to $k$---we will be thinking of
values of $M$ on the order of magnitude of $10^{18}$ and
$k<100$---then $(M-k)/M$ 
%
\mar{$(M-k)/M$ is essentially equal to 1}
%
will essentially equal~1, Hence
$$
\left( 1-\dfrac{\lambda}{M}\right)^{M-k}
   \approx \left(e^{-\lambda}\right)^1=e^{-\lambda}.$$

We can now work out $p_1$. Fix your attention first on a particular
atom.  The probability that {\em that\/} atom does emit an
$\alpha$-ray detected by the scintillation counter while the other
$M-1$ atoms do not is (again by Rule 3)
$$
\left(\dfrac{\lambda}{M}\right)^1
   \left( 1-\dfrac{\lambda}{M}\right)^{M-1}
   \approx \dfrac{\lambda}{M}\,e^{-\lambda},
$$
by the preceding approximation.

Since there are altogether $M$ atoms which might have been responsible
for the single $\alpha$-ray emission, the probability that some {\em
  unspecified\/} atom emits an $\alpha$-ray while the others do not is
(Rule 1, page~\pageref{rule1}) the sum of the probability we just
calculated for each of the $M$ atoms, which is equal to $M$ times that
probability.  The total probability is $p_1$:
$$
p_1 \approx M\,\dfrac{\lambda}{M}\,e^{-\lambda} 
= \lambda \,e^{-\lambda}.
$$

To work out $p_2$, note that the probability that each atom of some
fixed pair of atoms emits an $\alpha$-ray detected by the counter, and
no other atoms does, is
$$
\left( \dfrac{\lambda}{M} \right) \left( \dfrac{\lambda}{M} \right)
  \left(1- \dfrac{\lambda}{M}\right)^{M-2}
  \approx \dfrac{\lambda ^2}{M^2}\, e^{-\lambda},
$$
using our usual approximation.  Since there are $\frac{1}{2}M(M-1)$
different pairs of atoms (we can choose the first $M$ different ways
and the second $(M-1)$ ways, but each pair gets counted twice in this
scheme, so we have to divide by 2), we obtain
\[
p_2 \approx  \dfrac{M(M-1)}{2} \dfrac{\lambda ^2}{M^2}e^{-\lambda}
  =  \left( 1 - \dfrac{1}{M}\right) \dfrac{\lambda ^2}{2}e^{-\lambda}
  \approx \dfrac{\lambda ^2}{2}e^{-\lambda}.
\]

As one can easily imagine, the computations for $p_3$, $p_4$, $\ldots$
are similar.  The observant reader will note that the exact values we
got for $p_0$, $p_1$ and $p_2$ are not the values given by the Poisson
distribution.  We got the Poisson probabilities only by making various
approximations that were justified by the large value of $M$.  The
assumptions we have made actually lead to what is called the {\bf
  binomial distribution} (see chapter~12.1), a distribution which
tends to the Poisson distribution in the limit $M\rightarrow\infty$.
In this case, where $M$ is large and $\lambda$ relatively small, the
binomial distribution is extremely close to the Poisson distribution.
 
\subsubsection{Other applications of the Poisson distribution}

        The Poisson distribution can be used to model many other
situations that have a random element.  Examples include:
\begin{itemize}
\item The number of chromosome interchanges caused by exposure to
X-rays for a fixed interval of time.
\item The number of bacteria in a given unit of area on a Petri dish.
\item The number of misprints on a page in a book.
\item The number of flying-bomb hits per unit area in London during
World War II.
\end{itemize}
In the exercises we will explore some examples.



\subsection{Exercises}

\subsubsection{Probability models}
\index{probability model}

\vspace{-10pt}

\num A fair coin is tossed.  If it comes up $H$ (heads), a fair die is
rolled.  If the coin comes up $T$, the coin is tossed again.
Construct a probability model for this experiment, listing the
possible outcomes and their probabilities.  (Hint: the list of
outcomes is $H,1$, $H, 2$, $\ldots$, $H,6$, $TT$, $TH$.)


\num Two identical fair coins are put in cup, shaken, and spilled out
onto a table.  Construct a probability model for this experiment.


\numa In the disintegration of large numbers of particles of radium
($Ra$), it is noted that 29\% of the disintegrations result in
$$
Ra \longrightarrow P + A
$$ 
and the remainder in 
$$
Ra \longrightarrow He^+ + B.
$$
What is a model for the disintegration of a single particle of $Ra$?

\sub Construct a probability model for the disintegration of two
particles of~$Ra$.


\subsubsection{The Poisson distribution}

\vspace{-10pt}

\num The purpose of this exercise is to present another way to show
that the expected value $E$ of the Poisson distribution is equal to
$\lambda$.  As in the text we have
\[
E = 0 \cdot p_0 + 1 \cdot p_1 + 2 \cdot p_2 + 3 \cdot p_3
        + \cdots = \sum_{n=0}^\infty n p_n. 
\]
The numbers $np_n$ can be simplified as
follows:
\begin{align*}
0\cdot p_0 &=  0 ,\\[2pt]
1\cdot p_1 &=  1\cdot\lambda e^{-\lambda}  =
         \lambda\cdot e^{-\lambda}  =  \lambda p_0, \\[5pt]
2\cdot p_2 &=  2\cdot \dfrac{\lambda^2e^{-\lambda}}{2}  =
         \lambda\cdot\lambda e^{-\lambda}  =  
		\lambda p_1,  \\[10pt]
3\cdot p_3 &=  3\cdot \dfrac{\lambda^3e^{-\lambda}}{6}  =
         \lambda \cdot \dfrac{\lambda^2e^{-\lambda}}{2}  =
         \lambda p_2.
\end{align*}

\sub  This pattern generalizes: show that
\[np_n = \lambda p_{n-1}\quad\hbox{for all $n > 0$}\ .\]

\sub  Use part (a) to compute the expectation $E$ (you will need to 
use the fact that the sum of the
probabilities is $p_0+p_1+p_2+\cdots = 1$).


\num  A model is to be constructed for the number of rain drops that 
fall per square foot over a short time interval.  Under what conditions 
would a Poisson distribution be appropriate.  Under what conditions would 
a linear model be better?

\num In analyzing flying-bomb hits in the south of London during World
War II, investigators partitioned the are into 576 small sectors, each
being $\frac{1}{4}$ of a square kilometer.  There were 229 sectors
with no hits, 211 sectors with exactly 1 hit, 93 sectors with exactly
2 hits, 35 sectors with 3 hits, 7 sectors with 4 hits, and one sector
with 5 or more hits.  What might lead you to expect that a Poisson
distribution might be a good model for the number of hits on each
sector?  Fit a Poisson distribution to the data by taking $\lambda$ to
be the average number of hits per sector.  Use this $\lambda$ to
compute the theoretical frequencies of 0, 1, 2, 3, 4 and 5 hits in 576
sectors.

\num A meteorite shower sprinkles a large area of the earth's surface
with small meteorite hits.  The average density is $5\times 10^{-6}$
hits per square meter.  Set up a model assigning a probability to the
number of hits per square kilometer.

\num The central processing unit (CPU) of a laptop computer will
freeze if more than ten instructions are received in a millisecond.
If the average number of instructions per second received in the
course of executing a large program is one per millisecond, what is
the probability that the instructions received by the CPU will cause
it to freeze (and, hence, the program to crash).


\newpage
%%%%%%%%%%%%%%%%%%%%%%%%%%%%%%%%%%%%%%%%%%%%%%%%%%%%%%%%%%%%%%%%%%%%%%%
%                                                                     %
%                              Section 3                              %
%                                                                     %
%%%%%%%%%%%%%%%%%%%%%%%%%%%%%%%%%%%%%%%%%%%%%%%%%%%%%%%%%%%%%%%%%%%%%%%


\section{The Power Spectrum}


This section is an 
%
\mar{The problem of\\ signal + noise}	
%
application of ideas about periodic functions and integrals to the
problem of separating a signal from noise.  We face this problem in
our daily life.  Radio and television signals have noise added to them
from other radio sources we can't control.  The noise sounds like
hissing static on a radio and looks like ``snow'' on a television
screen.  A good receiver is designed to filter out the noise while
allowing the the transmitted signal to come through undistorted.
\vspace{150pt}

\special{psfile=../ps/calc2/lynx2.ps}

Scientific data and a radio broadcast have something in common: both
are combinations of signal and noise.  For instance, consider the
annual harvest of lynx pelts\index{lynx} by the Hudson's Bay
Company.\index{Hudson's Bay Company} It is conceivable that the lynx
population itself (the {\em signal\/}) was periodic, but various
random fluctuations (the {\em noise\/})\index{signal and noise} caused
the harvest (which is {\em signal\/} + {\em noise\/}) to take the form
it did.  \mar{The power spectrum filters noise to detect periodic
  signals} If this is the case, then we should try to ``filter out''
the noise and find the underlying periodic signal.  There is a
mathematical tool to do this; it is called the {\bf power
  spectrum}.\index{power spectrum} We will discuss the ideas behind
the power spectrum and show how it can be used to detect the
underlying in noisy data.

\subsection{Signal + Noise}

To prepare for working with the power spectrum, let's first see what
happens to a periodic signal that has some noise added to it.  The
signal we will use is a pure sine wave.  The {\bf
  information}\index{information} that the signal carries is the
frequency of that wave.  The noise will also be a function, but one
whose values vary in a random fashion.  It can be thought of as a
combination of periodic signals of all frequencies.  For this reason
it is sometimes called ``white noise,'' because white light is a
combination of light rays of all colors (i.e., frequencies).  Here is
the question we will explore: If we increase the strength of the
noise, when do we lose the information contained in the original
signal?
	
The signal and noise are shown below.  
%
\mar{A signal with\\ faint noise}	
%
As you can see, the amplitude of the signal is about 4 times as large
as the amplitude of the noise.  We say\linebreak

\vspace{16pt}

\special{psfile=../ps/calc2/noise2.ps}

\vspace{100pt}

\noindent
that the {\bf signal-to-noise ratio} is
4:1.\index{signal-to-noise ratio}
\label{faint noise}  
The combined signal + noise is no longer a pure sine wave, of course.
However, it is still recognizable as a ``noisy'' wave with the same
frequency as the original signal.  The information from the signal has
not yet been lost.

Look what happens when 
%
\mar{The noise level becomes stronger}	
%
we increase the amplitude of the noise.  In the figure below, the
noise has been increased by a factor of 4, so the signal-to-noise
ratio is now 1:1.  The combined signal + noise is now very noisy.
Would you be willing to argue that it is a wave of the same frequency
as the original signal?  Or would you prefer to say that it has no
periodic pattern whatsoever?  It appears we are close to losing the
information from the original signal.

\vspace{30pt}

\special{psfile=../ps/calc2/noise3.ps}

%\vspace{160pt}

\newpage

If we increase the original noise level 
%
\mar{The noise level becomes overwhelming}	
%
by a factor of 10, we appear to lose the original signal altogether.
The signal-to-noise ratio is now 1:2.5, and the \linebreak  

\vspace{20pt}

\special{psfile=../ps/calc2/noise4.ps}


\vspace{265pt}

\noindent
signal + noise appears to be as random as the noise itself.  In spite
of appearances, the signal is still there, and it will be detected in
the power spectrum!


\subsection{Detecting the Frequency of a Signal}\index{frequency!of a signal}

Assume we have a signal that may be distorted by a lot of noise.  
%
\mar{Compare the signal\\ to a probe whose frequency can be varied} 	
%
We want to decide whether the signal has a periodic component; if it
does, we want to determine its frequency.  Our detector is based on
this simple idea: {\em Compare the signal to a test probe of known
  frequency; vary the frequency of the probe until there is a positive
  response}.  Of course, we still need to explain how the comparison
is made, and what constitutes a positive response.
	
Although the detector will work on a very noisy signal, like the one
above, we will understand it better if we first use it to analyze a
signal whose periodic nature is evident.  
%
\mar{The test probe}	
%
Let the signal $S(t)$ be a pure sine wave lying above the $t$-axis,
and suppose that $t$ is the time measured in seconds.  Our {\bf test
  probe}\index{test probe} is the function
	$$
	P(t) = \sin(2\pi \omega t)
	$$
whose frequency is $\omega$ cycles per second.  As its graph demonstrates, the values of $P$ are equally likely to be positive or negative.  

Is the same true for the product $P(t)S(t)$?  
%
\mar{When the signal matches the\\ test frequency}	
%
Suppose first that $S(t)$ has the same frequency as $P(t)$ (below,
left).  As you can see, the positive values of $P(t)$ are always
multiplied by the larger values of $S(t)$.  By contrast, the negative
values of $P(t)$ are always multiplied by the smaller values of
$S(t)$.  Consequently, the positive values of $P(t)S(t)$ outweigh the
negative ones.  On average, the value of the product is positive.  In
fact, the average value\index{average value} of the product is half
the amplitude of the original signal. \label{frequencies match} Later
on we will see why this is so.
\vspace{80pt}

\special{psfile=../ps/calc2/sintest1.ps}  	  

\vspace{195pt}

On the right we see what happens 
%
\mar{When the signal {\em doesn't\/} match the\\ test frequency}	
%
if $S(t)$ is {\em not\/} related to $P(t)$.  In that case, a large
value of $S(t)$ is just as likely to multiply a positive value of
$P(t)$ as a negative one.  Consequently, the product $P(t)S(t)$ will
have both large positive and large negative values.  On average, the
value of the product will be about 0.

Let's use the detector on the signals we constructed on
page~\pageref{faint noise}.  In both we started with a pure sine wave
and added some white noise.  In the first, the signal-to-noise ratio
was 4:1.
\vspace{60pt}

\special{psfile=../ps/calc2/sintest2.ps}
	
\vspace{160pt}	
	
\noindent
In the second the noise was stronger; the signal-to-noise ratio was 1:1.
\vspace{70pt}

\special{psfile=../ps/calc2/sintest3.ps}
	
\vspace{190pt}	
	
To use the detector yourself, you have to be able to calculate the
average value of a function.  This is discussed in chapter~6.3.  The
average value of $y = f(x)$ on the interval
%
\mar{\vspace{14pt} The average value\\ of a function}
%
$a \leq x \leq b$ is
\[
\dfrac{1}{b - a} \int_a^b f(x) \, dx.
\]
Our detector is the average value of the product of the signal $S(t)$
and the test probe $P(t) = \sin(2\pi \omega t)$.\index{frequency
  detector}\index{average value}
\begin{center}
\begin{picture}(340,50)
\put(0,0) {\framebox(340,50)
{
\begin{minipage}{310pt}
{\bf Frequency detector}: \hfill $D(\omega) = \dfrac{1}{b-a} \dint_a^b S(t) \sin(2\pi\omega t)\,dt$. \hfill \mbox{}
\end{minipage}
}}
\end{picture}
\end{center}

Clearly, the value of the detector depends 
%
\mar{Integrals\\ with parameters\\ define functions}
%
on the frequency $\omega$ of the probe $P$.  We have tried to reflect this in the notation: the detector is a function $D$ whose input is the frequency $\omega$.  The output of the function is calculated as an integral in which the input $\omega$ plays the role of a parameter.  

This is the first time we have defined a function as an integral with
a parameter.  Let's see how the detector works to analyze the signal
$S(t) = 3\sin(5t)$ over the interval $0 \leq t \leq 10$.  We have
	$$
D(\omega) = \dfrac{1}{10} \int_0^{10} 3\sin(5t) \sin(2\pi\omega t)\,dt.
	$$
In the exercises at the end of chapter~11.3 we obtained an explicit
formula for the integral of the product of two sine functions.  Find
that formula and check that it yields the following:
	$$
	D(\omega) = \dfrac{3}{10(4\pi^2\omega^2 - 25)}
	(5\cos(50)\sin(20\pi\omega)-2\pi\omega\sin(50)\cos(20\pi\omega)).
	$$	
Notice, in your own calculations, that $\omega$ emerges as the
variable on which the whole expression depends.

The graph of $D(\omega)$ is shown on the top of the next page.  You
should plot it yourself, using a computer graphing utility.
%
\mar{$D(\omega)$ peaks when\\ $\omega$ is the frequency\\ of the signal}
%
For most frequencies $\omega$, the value of the detector $D$ is close
to 0.  There is a single strong peak, which you can find at $\omega
\approx .795$ cycles/sec.  As it happens, the frequency of the signal
$S = 3\sin(5t)$ is $5/2\pi = .79577\ldots$ cycles/sec!  Moreover, the
height of the peak is about 1.5, which is exactly half the amplitude
of the signal.

\newpage

\mbox{}
\vspace{80pt}

\special{psfile=../ps/calc2/detect1.ps}

\vspace{50pt}

The graph above tests the signal $S(t)$ when the detector is
integrated over a time interval that is 10 seconds long.  That is, $0
\leq t \leq 10$ seconds.  If we repeat the test by integrating over a
much larger interval, the frequency detector gives us a sharper report
on the frequency of the signal.  In the graph below the function
$D(\omega)$ was calculated by integrating over the interval
%
\mar{\vspace{24pt} The peak in $D(\omega)$\\ is sharper if the signal\\ is
  tested over a\\ longer time interval} 
%
$0 \leq t \leq 100$ seconds.
\vspace{100pt}

\special{psfile=../ps/calc2/detect2.ps}

%NOTE: this won't run on a 300 dpi laser printer unless the stepsize in
%the graphing program is .3 or greater

\vspace{65pt}

\noindent
{\bf Computation}.  Of course, it is rare to find a formula for
$D(\omega)$ in terms of the frequency $\omega$.  For most signals
$S(t)$, the best we can do is calculate the value of the integral
numerically for a sequence of values of the parameter $\omega$.  The
program DETECTOR,
%
\mar{The program DETECTOR}	
%
which is listed on the next page, does this.  As it is written, it
analyzes the function $3\sin(5t)$ on the interval $0 \leq t \leq 10$,
and it produces the graph $D(\omega)$ at the top of this page.  The
``outer loop''
	\begin{quote}
	{\tt FOR j = 1 TO omegasteps} \quad \ldots \quad {\tt NEXT j}
	\end{quote}
plots $D(\omega)$ over the interval $0 \leq \omega \leq 3$, using
$2^{10}$ equally spaced values of $\omega$.  Each $D(\omega)$ is an
integral whose value is first calculated as a midpoint Riemann sum
with $2^7$ steps.  The calculation is carried out by the short ``inner
loop''
	\begin{quote}
	{\tt FOR k = 1 TO numberofsteps} \quad \ldots \quad {\tt NEXT k},
	\end{quote}
which you should recognize as an adaptation of the program RIEMANN
from chapter~6.


\begin{center}
{\bf Program: DETECTOR \\  \index{program!DETECTOR}
To detect the frequency of a signal}
\end{center}

\footnotesize
\hspace{-20pt}
\begin{minipage}{358pt}
\noindent
{\sl Set up GRAPHICS}
\vspace{-11pt}
\begin{verbatim}
startomega = 0
endomega = 3
omegasteps = 2 ^ 10
deltaomega = (endomega - startomega) / omegasteps
twopi = 8 * ATN(1)
DEF fnf (t) = 3 * SIN(5 * t)
a = 0
b = 10
numberofsteps = 2 ^ 7
deltat = (b - a) / numberofsteps
omega = startomega
oldomega = omega
oldaccum = 0
FOR j = 1 TO omegasteps
     t = a + deltat / 2
     accum = 0
     FOR k = 1 TO numberofsteps
          deltaS = (fnf(t) * SIN(twopi * omega * t) * deltat) / (b - a)
          accum = accum + deltaS
          t = t + deltat
     NEXT k
     omega = omega + deltaomega
\end{verbatim}

\vspace{-10pt}
\rule{23pt}{0pt} {\sl Plot the line from {\tt (oldomega, oldaccum)} to {\tt (omega, accum)}}

\vspace{-13pt}
\begin{verbatim}
     oldomega = omega
     oldaccum = accum
NEXT j
\end{verbatim}
\end{minipage}
\normalsize
\bigskip

If we modify the program DETECTOR so that it analyzes the function
\[
S(t) = 3\sin(5t) + \sin(8t),
\]
we get the graph at the top of the next page.  The scale on the
$\omega$-axis has also been modified to make it easier to read
multiples of $1/2\pi$ cycles per second.  Notice the strongest peak is
at $\omega = 5/2\pi$ cycles/sec, and $D \approx 1.5$ there.  But there
is now a second peak at $\omega = 8/2\pi$ cycles/sec, where $D \approx
.5$.  Indeed, $S$ consists of two periodic components, one with three
times the amplitude of the other.  The stronger component has
frequency $5/2\pi$ cycles/sec, the weaker $8/2\pi$ cycles/sec.
\vspace{95pt}

\special{psfile=../ps/calc2/detect3.ps}

\vspace{25pt}

The following example first appeared in chapter~7.2.  It is clear from
the graph that it has a basic frequency of 5 Hz.  The detector shows
that it also has an equally strong component at 10 Hz and a much
weaker component at 15 Hz.
%
\mar{\vspace{40pt} A periodic \\ signal \ldots \\ \vspace{100pt}
  \ldots and its\\ frequency detector}
%
Can you guess a formula for $g(t)$?
\vspace{75pt}

\special{psfile=../ps/calc2/period3.ps}
 
\vspace{140pt}

\special{psfile=../ps/calc2/detect4.ps}

\vspace{25pt}

\noindent
The graph of $z = D(\omega)$ was produced by DETECTOR.  The integral was calculated for {\tt a = 0}, {\tt b = 10}, and {\tt numberofsteps = 2 \verb+^+ 9}.    


\subsection{The Problem of Phase}\index{phase}

\noindent
\begin{minipage}[t]{300pt}
Our detector is built on the premise that, if you take the product of
two functions of the same frequency, its average value\index{average
  value} will be different from~0.  This is illustrated by the top
three graphs on the left.  The signal and the probe are both
$\sin(t)$.  Their product is a function that ranges between 0 and 1,
and has average value 1/2.  However, something quite different happens
if we change the signal from $\sin(t)$ to $\cos(t)$.  This doesn't
change the period, but it does change the product, as you can see in
the three lower graphs.  The new product is centered around the
$t$-axis; its average value is~0.  Thus the detector fails to reveal
that the signal has the same frequency as the probe.
		  
\quad \ A closer look at the two sets of graphs will show what has
happened.  In the first case, when $P$ is positive, so is $S$.  When
$P$ is negative, so is $S$.  Thus, the product $S \cdot P$ is never
negative; on average, its value is positive.  This is what we expect.

\quad \ The second case is only a little more complicated.  When $P$
is positive, $S$ is positive only half the time; the other half it is
negative.  Consequently, the product $S \cdot P$ takes both positive
and negative values.  The same thing happens when $P$ is negative.  On
average, the value of the product is~0, {\em even though the
  frequencies of $P$ and $S$ match}.
	
\quad \ The problem is that their {\em phases\/} don't match.  The
signal $S = \cos(t)$ hits its peak $\pi/2$ seconds before the probe $P
= \sin(t)$.  This kind of a difference is called a {\bf phase shift}.
In the exercises for chapter~7.2, you showed that if the phase of the
sine function is shifted to the left by $\pi/2$, the result is the
cosine function:
\vspace{-8pt}
	$$
	S = \sin(t + \pi/2) = \cos(t).
	$$
\end{minipage}
\hfill
\begin{minipage}[t]{40pt}
\special{psfile=../ps/calc2/sintest6.ps}

\vspace{220pt}
\special{psfile=../ps/calc2/sintest5.ps}

\end{minipage}





\medskip

\noindent
Since $\pi/2$ radians is the same as 90$^\circ$, we sometimes express
this equation by saying that ``the sine and the cosine are 90$^\circ$
out of phase.''
	
Of course the signal could involve a phase shift 
%
\mar{Arbitrary phase shifts}\index{phase!shift}	
%
of any amount $\varphi$: $S = \sin(t - \varphi)$.  All these signals
have the same period as the probe $P = \sin(t)$.  Exercise 20 of
chapter 7.2 shows what happens if this signal is tested against the
probe: the average value of the product $S \cdot P$ is
$\cos(\varphi)/2$.  Clearly, this depends on the size of the phase
shift $\varphi$.  In particular, if \mar{The average value varies with
  the phase} $\varphi = 0$ (so $S = \sin(t)$), the average value
is~1/2.  If $\varphi = -\pi/2$ (so $S = \cos(t)$), the average value
is~0.  The formula therefore agrees with what we already know for the
two signals we considered as examples.

There is one more case worth glancing at: $\varphi = \pm \pi$.  This
is also called a phase shift of 180$^\circ$.  It doesn't matter
whether you go forward 180$^\circ$ or backward; in either case $S =
\sin(t \pm \pi) = -\sin(t)$.  This time the average value of the
product is~$-1/2$.

The problem of phase 
%
\mar{The problem\\ of phase \ldots}	
%
is now be clear: The probes $P = \sin(2 \pi \omega t)$ have trouble
detecting the frequency of a signal that is out of phase with them.
However, any phase-shifted sine function can be expressed as a sum of
pure sine and cosine functions:
\[
\sin(bt - \varphi) = M\sin(bt) + N\cos(bt), 
\] 
where $M = \cos(\varphi)$ and $N = -\sin(\varphi)$.  (See the
exercises.)  Since the sine probes $P$ will detect $M\sin(bt)$,
%
\mar{\ldots and its solution}	
%
we need only construct a second set of probes to detect $N\cos(bt)$.
The test probes we add are the cosine functions
	$$
	P_c = \cos(2 \pi \omega t).
	$$ 
We use the subscript ``$c$'' to distinguish these from the sine
probes, which henceforth will be denoted~$P_s$.

We must also construct a second detector, 
%
\mar{Two new detectors}	
%
to handle the new cosine probes.  Let's take this opportunity to make
a technical adjustment: we redefine a detector to be {\em twice\/} the
average value of the signal and the probe.  In that way, the height of
the detector at a peak equals the amplitude of the signal at that
frequency---rather than half the amplitude. \label{detectors}
\index{sine detector}\index{cosine detector}

\begin{center}
\begin{picture}(340,85)
\put(0,0) {\framebox(340,85)
{
\begin{minipage}{310pt}
\rule{14pt}{0pt}
	{\bf Sine detector}: 
	\hfill $D_s(\omega) = 
	\dfrac{2}{b-a} \dint_a^b S(t) \sin(2\pi\omega t)\,dt$. 
	\hfill \mbox{}\\[8pt]
{\bf Cosine detector}: 
	\hfill $D_c(\omega) = 
	\dfrac{2}{b-a} \int_a^b S(t) \cos(2\pi\omega t)\,dt$. 	
	\hfill \mbox{}
\end{minipage}
}}
\end{picture}
\end{center}

You can modify the program DETECTOR 
%
\mar{The graphs of\\ $D_s$ and $D_c$}
%
to produce the graphs of $D_s(\omega)$ and $D_c(\omega)$.  You can see
below how they analyze the signal $S = \cos(7t)$ over the interval $0
\leq t \leq 10$.  The cosine detector $D_c$ has a shape we've seen
\linebreak


\mbox{}

\vspace{60pt}

\special{psfile=../ps/calc2/detect5.ps}
	
\vspace{170pt}

\noindent
before.\label{cos7t} It has a single peak at $\omega = 7/2\pi$
cycles/sec, which is the frequency of the signal.  The peak is 1 unit
high, which is the amplitude of the signal.  The sine detector has an
unfamiliar shape.  Notice first that $D_s(7/2\pi) = 0$.  This confirms
our earlier observation that the average value of the product of a
sine and a cosine at the same frequency is~0.  For values of $\omega$
slightly larger or smaller than $7/2\pi$, though, the sine detector
swings relatively far from~0.  This pattern is typical when a detector
is analyzing a signal that is 90$^\circ$ out of phase with the probes.
\medskip

\noindent
{\bf Resonance}.\index{resonance} Try this experiment.  Sit at a piano
and hold all the pedals down.  Then sing a note.  If you sing loud
enough, and hold the note long enough, one of the piano strings will
start vibrating.  If you stop abruptly and listen to the string, you
will hear it sounding the same note you were singing.
%
\mar{The physical analogue of a detector\\ is a resonator}
%
The piano has detected the frequency of your signal!  It is the
physical analogue of our mathematical frequency detectors.  The
response of the string is called {\bf resonance}.  Had you sung a
lower note, a larger string would have resonated.

Resonance gives us a vivid language for describing how our detectors
work.  We can say a test probe ``resonates'' with a signal when their
product is different from zero on average.  The larger the average
value, the stronger the resonance.

\newpage

\aside{Resonance occurs all around us.  Sometimes it is a
  nuisance---for instance, when the windows in our house rattle while
  a heavy truck drives by, or an air conditioner runs.  Sometimes we
  exploit it deliberately---for instance, when we use a radio tuner as
  an electronic resonator to detect and amplify certain
  electromagnetic waves.}


\subsubsection{Detector as transform}\index{transform}

We now have two distinct ways to describe a signal $S$.  The function
$S(t)$ is one way.  It tells us how strong the signal is at each
instant $t$.  But we can also think of the signal as a mixture of sine
and cosine waves of different frequencies.  The detectors
$D_s(\omega)$ and $D_c(\omega)$ tell us how strong the signal is at
each frequency $\omega$.  That is the second way.
	
There is a direct connection between these two descriptions, of course.  It is provided by the formulas
	$$
	D_s(\omega) = \dfrac{2}{b-a}
			\int_a^b S(t) \sin(2\pi\omega t)\,dt
		\qquad
	D_c(\omega) = \dfrac{2}{b-a}
			\int_a^b S(t) \cos(2\pi\omega t)\,dt.
	$$
In effect, these formulas tell us how to {\em transform\/}  
%
\mar{Integrals transform\\ $S$ into $D_s$ and $D_c$}
%
the first description $S(t)$ into the second $D_s(\omega)$,
$D_c(\omega)$.  The transformation is so complete that even the input
variable is changed---from $t$ to $\omega$.  Look back at the formulas
to see how the new variable $\omega$ is brought in.

Our detectors are essentially the same as the {\bf Fourier sine
  transform}\index{Fourier transform} and the {\bf Fourier cosine
  transform}.  There is also an {\bf inverse Fourier
  transform}\index{Fourier transform!inverse} that works in reverse:
it produces $S(t)$ from the frequency data $D_s(\omega)$ and
$D_c(\omega)$.  The Fourier transforms are an important tool in
mathematics and in science.  For example, a hologram is the Fourier
transform of an ordinary image.  Fourier transforms and their inverses
are used in photo restoration, in the enhancement of the digitized
pictures sent back from cameras in space, and in filtering the signal
in a stereo set.
	
\aside{The French mathematician Jean Baptiste Fourier (1768--1830)
  introduced what we call Fourier transforms and Fourier series to
  study the conduction of heat.  Now his methods are used to study all
  sorts of periodic and non-periodic phenomena.  They are also the
  foundation for the part of pure mathematics called harmonic
  analysis.}\index{Fourier, J. B.}


\subsection{The Power Spectrum}

The sine and cosine detectors 
%
\mar{A detector that ignores phase differences}\index{phase}	
%
provide enough information to reconstruct the original signal in
complete detail---including phase.  Often, though, they provide more
detail than we want.  
\pagebreak
We can use another tool---called the {\bf power
  spectrum}---to\index{power spectrum} determine only the strength of
the different frequencies that occur in a signal, without regard to
their phase. The power spectrum is constructed from the two detectors
in the following way:
\begin{center}
\begin{picture}(330,40)
\put(0,0) {\framebox(330,40)
{
\begin{minipage}{300pt}
{\bf Power spectrum}: \hfill $P(\omega) = \displaystyle{\sqrt{[D_s(\omega)]^2 + [D_c(\omega)]^2}}$ \hfill \mbox{}
\end{minipage}
}}
\end{picture}
\end{center}
  
To see how the power spectrum works, we'll consider the signal $S(t) =
A\sin(7t - \varphi)$.  This is a sine wave of frequency $\omega =
7/2\pi$ and amplitude $A$.  Let's concentrate first on $\omega =
7/2\pi$.  If there were no phase shift $\varphi$ present, we would
expect that
	$$
	D_s(7/2\pi) = A \qquad D_c(7/2\pi) = 0.
	$$
However, because there is a phase shift, the actual values turn out to be
$$
D_s(7/2\pi) = A\cos{\varphi} \qquad D_c(7/2\pi) = -A\sin{\varphi}.
$$
(These calculations are given as exercises.)  The values of the
detectors clearly depend on the phase shift.  By contrast,
\begin{align*}
P(7/2\pi) &= \displaystyle{\sqrt{[D_s(7/2\pi)]^2 + [D_c(7/2\pi)]^2}}\\
	  &= \displaystyle{\sqrt{A^2\cos^2{\varphi} + 
			A^2\sin^2{\varphi}}}\\
	  &= A.
\end{align*}
We have used the fact that $\cos^2{\varphi} + \sin^2{\varphi} = 1$ for
every $\varphi$.  Thus, the power spectrum does {\em not\/} depend on
the phase.  It tells us only the amplitude of the signal at the
frequency $\omega = 7/2\pi$.

If we calculate the power spectrum over all frequencies $\omega$, we
get the graph shown at the top of the next page.
%
\mar{The program POWER}	
%
The program POWER generates this graph. It was derived from the
program DETECTOR.  Compare the two programs, particularly the terms
{\tt deltaS} and {\tt deltaC}.  In POWER, they have been multiplied
by~2, to agree with our new definition of $D_s$ and $D_c$ on
page~\pageref{detectors}.
	
\mar{\vspace{25pt} \mbox{Power spectrum of} $3\sin(7t - \pi/3)$}
\vspace{80pt}

\special{psfile=../ps/calc2/detect6.ps}

\vspace{20pt}

\begin{center}
{\bf Program: POWER \\  \index{program!POWER}
The power spectrum of a signal}
\end{center}

\footnotesize
\hspace{-20pt}
\begin{minipage}{358pt}
\noindent
{\sl Set up GRAPHICS}
\vspace{-11pt}
\begin{verbatim}
startomega = 0
endomega = 3
omegasteps = 2 ^ 9
deltaomega = (endomega - startomega) / omegasteps
pi = 4 * ATN(1)
twopi = 2 * pi
DEF fnf (t) = 3 * SIN(7 * t - pi / 3)
a = 0
b = 10
numberofsteps = 2 ^ 6
deltat = (b - a) / numberofsteps
omega = startomega
oldomega = omega
oldpower = 0
FOR j = 1 TO omegasteps
     t = a + deltat / 2
     accumS = 0
     accumC = 0
     power = 0
     FOR k = 1 TO numberofsteps
          deltaS = 2 * (fnf(t) * SIN(twopi * omega * t) * deltat) / (b - a)
          accumS = accumS + deltaS
          deltaC = 2 * (fnf(t) * COS(twopi * omega * t) * deltat) / (b - a)
          accumC = accumC + deltaC
          t = t + deltat
     NEXT k
     power = SQR(accumS ^ 2 + accumC ^ 2)
     omega = omega + deltaomega
\end{verbatim}

\vspace{-10pt}
\rule{23pt}{0pt} {\sl Plot the line from {\tt (oldomega, oldpower)} to {\tt (omega, power)}}

\vspace{-13pt}
\begin{verbatim}
     oldomega = omega
     oldpower = power
NEXT j
\end{verbatim}
\end{minipage}
\normalsize

\medskip

To see how the power spectrum 
%
\mar{Two signals whose components differ\\ only in phase}
%
detects the frequencies in a signal while overlooking the phases of
the different components, consider these two
signals:\label{testfunction}
	\begin{eqnarray*}
	g(t) &=& 10\sin(7t) + 7\cos(13t) + 5\cos(23t) \\
	h(t) &=& 10\sin(7t) + 7\cos(13t) - 5\cos(23t)
	\end{eqnarray*}
They differ only in the sign of the last term.  This is equivalent to
a phase shift of 180$^\circ$ in that term.  The graphs are drawn below
(with constants \linebreak

\mbox{}

\vspace{60pt}

\special{psfile=../ps/calc2/period4.ps}

\vspace{20pt}
\noindent
added to separate them vertically).  It is remarkable how
different the graphs appear to be, considering how nearly alike their
formulas are.  You can find similarities if you look closely, though.
For instance, the peaks of one graph tend to match the peaks of the
other.

\vspace{235pt}

\special{psfile=../ps/calc2/detect7.ps}

The power spectrum, however, has no trouble detecting the similarities
between the two signals.  As you can see, they indicate that the same
dominant frequencies occur in $g$ and $h$, and that corresponding
frequencies occur with the same amplitude.  We learn that the formula
for $g$ or $h$ can be written as
	$$
	10\sin(7t - \varphi_1) + 7\sin(13t - \varphi_2) + 
		5\sin(23t - \varphi_3).
	$$
The only thing we can't learn from the power spectrum are the three
phase differences $\varphi_1$, $\varphi_2$, $\varphi_3$.

\newpage

The graphs of the power spectra were drawn by POWER, using the
following values:
	\begin{verbatim} 
     endomega = 4
     omegasteps = 2 ^ 8
     numberofsteps = 2 ^ 7
	\end{verbatim}
These two graphs actually differ very slightly.  You can see the
difference most clearly near $\omega = 20/2\pi$.
\medskip

For a final demonstration of the properties of the 
%
\mar{Detecting a periodic wave in a noisy signal}	
%
power spectrum, we return to the signal + noise problem that we raised
at the beginning of this section.  Let's see what happens to the power
spectrum of a pure sine wave when we gradually gradually add noise.
For simplicity, we take the frequency of the pure signal to be 2
cycles/sec.  The spectrum has a single strong spike at this frequency.
\vspace{185pt}

\special{psfile=../ps/calc2/detect8.ps}

\vspace{25pt}

One the following pages you can see what happens as the noise level is
increased. The power spectrum, which was virtually zero for all
$\omega > 3$ cycles/sec, is now non-zero for almost all frequencies in
the range we have graphed.
%	
\mar{In the power spectrum, noise and signal are separated}	
%
In other words, the noise is a mixture of many frequencies.  Notice
how the height of the power graph increases with the strength of the
noise.  This is most noticeable in the higher frequencies.
Eventually, in the final graph, we lose sight of the signal; the noise
has swamped it.  The signal to noise ratio is 1:2.5, meaning that the
noise is 2$\frac{1}{2}$ times as strong as the signal.  Nevertheless,
the power spectrum still shows a strong spike at $\omega = 2$
cycles/sec.  This corresponds to the signal.  The power spectrum can
still see the signal even when we can't!

\vspace{180pt}

\special{psfile=../ps/calc2/detect9.ps}

\vspace{230pt}

\special{psfile=../ps/calc2/detect10.ps}

\newpage

\mbox{}
\vspace{235pt}

\special{psfile=../ps/calc2/detect11.ps}

\vspace{25pt}

\subsection{Exercises}

\subsubsection{The problem of phase}

\vspace{-10pt}

\num  Use the ``sum of two angles formula,''
	$$
	\sin(A + B) = \sin(A) \cos(B) + \cos(A) \sin(B),
	$$
to show that the circular function $\sin(bt - \varphi)$ with period
$2\pi/b$ and phase difference $\varphi$ can be written as a
combination of pure sine and cosine functions of the same period:
	$$
	\sin(bt - \varphi) = M \sin(bt) + N \cos(bt).
	$$
show that $M = \cos(\varphi)$ and $N = -\sin(\varphi)$.  [Note that
  $M^2 + N^2=1$.]

\numa Express $\sin(5t - \pi/3)$ as a sum of a pure sine function and
a pure cosine function.

\sub Express $\frac{\sqrt{3}}{2}\sin(7t) + \frac{1}{2}\cos(7t)$ in the
form $A\sin(bt - \varphi)$.  To check your result, graph it together
with the given function using a computer graphing utility.

\sub Express $f(t) = \sin(t) + 2\cos(t)$ in the form $A\sin(bt -
\varphi)$.  Notice that the formula in exercise 1 requires that $M^2 +
N^2 = 1$, but in this example $M^2 + N^2 = 5$.  Therefore, first write
	$$
	f(t) = \sqrt{5} \left( \tfrac{1}{\sqrt{5}} \sin(t) + 
		\tfrac{2}{\sqrt{5}} \cos(t) \right).
	$$
The expression in parentheses has the right form.  Does your result
check on a computer?

\num  Suppose 
	$$
	A \sin(bt - \varphi) = M \sin(bt) + N \cos(bt).
	$$
How are $A$, $M$, and $N$ related?

\num The functions $\sin(t) + 2 \cos(t)$ and $2\sin(t) + \cos(t)$ have
the same period but differ in phase.  What is the phase difference?
Determine this two ways: by graphing, and by writing each expression
as a single function of the form $A \sin(bt - \varphi)$.

\num  Choose values for $A$, $b$, and $\varphi$ so that the function
	$$
	3 \sin(2x) + 4 \cos(2x) + A \sin(bx - \varphi) 
	$$
is {\em identically\/} zero---that is, equal to 0 for every value of~$x$.

\num  Choose values of $A$ and $\varphi$ so that the function
	$$
	\sin(x) + \sin(x + 1) + \sin(x + 2) + A\sin(x - \varphi)
	$$
is identically zero.


\subsubsection{The programs DETECTOR and POWER}


The purpose of these exercises is to give you experience interpreting
the power spectrum of a known signal using the program POWER and
modifications of DETECTOR.  The first exercise asks you to construct
these modifications.
	
\num Modify DETECTOR to produce two new programs, SDETECTOR and
CDETECTOR, which generate the sine detector and the cosine detector
functions that appear on page~\pageref{detectors}.

\numa Compare the outputs of $f(t) = \sin(t)$ and $g(t) = \cos(t)$ on
POWER.  Use the domain $0 \leq \omega \leq 1$.  Does POWER distinguish
between these functions?  Would you expect it to?

\sub Compare $f(t)$ and $g(t)$ using SDETECTOR.  Does SDETECTOR
distinguish between these functions?  Would you expect it to?

\sub Compare $f(t)$ and $g(t)$ using CDETECTOR.  Does CDETECTOR
distinguish between these functions?  Is the output of $g(t)$ on
CDETECTOR the same as the output of $f(t)$ on SDETECTOR?

\numa Describe the power spectrum of the signal $S = \sin(t) +
\cos(t)$.  How many peaks are there, and where are they?

\sub How does the spectrum of $S$ compare with the two generated in
the last question?

\sub Describe the output of SDETECTOR and CDETECTOR for the signal
$S$.  Compare these outputs to the corresponding outputs for $f$ and
$g$ in the last exercise.

\numa   Graph the function
	$$
	h(t) = 10\sin(7t) + 7 \cos(13(t) - 5 \cos(23t)
	$$
over the domain $0 \leq t \leq 14$.  Compare your result with the
graph on page~\pageref{testfunction}.

\sub Graph the power spectrum of $h(t)$ over the frequency domain $0
\leq \omega \leq 4$.  Compare your result with the text.  How many
peaks are there?  Where are they?  How high are they?  Do these
results agree with the amplitude and frequency information provided by
the formula for $h(t)$?

\num (Continuation of the previous exercise.)  Use SDETECTOR to
analyze $h(t)$ over the same frequency domain.  Compare the pattern
near $\omega = 13/2\pi$ with the patterns generated by the sine and
cosine detectors that appear on page~\pageref{cos7t}.  Compare the
patterns near $\omega = 7/2\pi$ and near $\omega = 23/2\pi$ the same
way.  Would you expect the patterns near $\omega = 13/2\pi$ and
$\omega = 23/2\pi$ to be similar?  Are they?  Are they similar to the
pattern near $\omega = 7/2\pi$?  Is this what you would expect?

\num (Continuation.)  Use CDETECTOR to analyze $h(t)$.  Follow the
guidelines of the previous question.


\subsubsection{A Grain of Salt}

The purpose of the power spectrum is to make visible the periodic
patterns contained with a given function.  However, our {\em method of
  computing\/} the spectrum can introduce spurious information, too.
It can tell us there are periods that are not really present in the
function.  So we must take the calculations with a grain of salt.  The
purpose of these exercises is to point out the spurious information,
show why it arises, and how we can get rid of it.


\num Use the program POWER to graph the power spectrum of the function
$\sin{2 \pi x}$ on the interval $0 \leq x \leq 10$.  Let $0 \leq
\omega \leq 3$.  Set
\begin{verbatim}
     numberofsteps = 100
\end{verbatim}
but let all the other parameters keep the values they have in the
program.

\ans{The power spectrum has a single peak of height~1 at $\omega
  \approx 1$}


\num Now increase the domain of integration to $0 \leq x \leq 30$, and
set 
\begin{verbatim}
     numberofsteps = 300
\end{verbatim}
to adjust for the increase in the size of the domain.  Use POWER again
to graph the power spectrum.  Compare this spectrum with the previous
one.

\num Leave $0 \leq x \leq 30$, but restore {\tt numberofsteps = 100}.
Use POWER once again to graph the power spectrum.  Compare this
spectrum with the previous two.

\ans{A new peak, of height 1, appears at $\omega \approx 7/3$.}

\num Let {\tt numberofsteps = 50}, and calculate the power spectrum
one more time.  What happens?
\medskip

When we reduce the number of integration steps, new peaks appear in
the power spectrum.  These new peaks represent {\em spurious\/}
information: the function $\sin{2\pi x}$ has no components whose
frequencies are 2/3, 7/3, or 8/3.  Let's see why this happens.  We'll
concentrate on $\omega = 7/3$.  First, you must decide whether the
peak in the power spectrum at $\omega = 7/3$ comes from the sine or
the cosine detector.

\num Use SDETECTOR and CDETECTOR to analyze $\sin(2\pi x)$.  Take $0
\leq x \leq 30$, $0 \leq \omega \leq 3$, and set {\tt numberofsteps =
  100}.  One of these detectors has the value 0 when $\omega = 7/3$.
Which one?


\num According to the previous exercise, the peak in the power
spectrum that is detected at $\omega \approx 7/3$ comes from the
integral
	$$
	\dfrac{2}{30}\int_0^{30} \sin(2\pi x) 
	\sin\left(2\pi \tfrac{7}{3} x\right)\, dx,
	$$
{\em not\/} from the cosine integral.  By using one of the sine and
cosine integrals from the exercises for chapter~11.3, determine the
{\em exact\/} value of this integral.  Is this the value you expected
to get?
\medskip

The program POWER calculates the spectrum numerically.  In particular,
we used it to calculate
	$$
	\int_0^{30} \sin(2\pi x) 
	\sin\left(2\pi \tfrac{7}{3} x\right)\, dx,
	$$
with 100 steps.  The step size is therefore $\Delta x = .3$.  In the
following exercises you will duplicate this numerical work ``by
hand.''


\num  Make a sketch of the graph of the function
	$$
 h(x) = \sin(2\pi x) \sin\left(2\pi \tfrac{7}{3} x \right)
	$$
on an appropriate interval.  What is the period of this function?

\num Determine the value of $h(x)$ at $x = 0$, .3, .6, .9, 1.2, and
1.5, and use these values to construct a Riemann sum for the integral
	$$
		\int_0^{1.5} h(x) \, dx
	$$
using left endpoints and a step size of $\Delta x = .3$.  Mark these
values of $h$ on the sketch you made in the previous exercise.

\ans{The Riemann sum is $-.3(2 \sin^2(2\pi/5) + 2 \sin^2(\pi/5)) = -.75$.}


\num  Evaluate the expression 
$$
\dfrac{1}{15}\int_0^{30} h(x) \, dx
$$
using a left endpoint Riemann sum with a step size of $\Delta x = .3$
How can you use the previous exercise to answer this question?

\ans{$-1$.  Since $h(x)$ is periodic with period $x = 1.5$, the
  interval $[0, 30]$ contains 20 periods of $h$.  The integral of $h$
  over $[0, 30]$ is therefore 20 times its integral over $[0, 1.5]$.}

\num  Compare the values of the detector
	$$
	\dfrac{2}{30}\int_0^{30} \sin(2\pi x) 
	\sin\left(2\pi \tfrac{7}{3} x\right)\, dx,
	$$
you have obtained by antidifferentiation and by numerical integration.

\medskip

These exercises demonstrate that the exact and computed values of the
power spectrum can be quite different, essentially because the steps
in a Riemann sum can pick out very special values of the integrand.
	
One way to deal with the problem is 
%
\mar{The {\em true\/} spectrum\\ is the limit of the computed graphs\\ of the spectrum}	
%
to increase the number of steps.  How will you know if you have gone
far enough?  Increase in stages until the graph of the power spectrum
{\bf stabilizes}---that is, until it no longer changes when you make a
further increase in the number of steps.

Of course, increasing the number of steps increases computer time.
This creates new problems.  To deal with them, however, we can switch
to more efficient numerical integration methods.  Simpson's rule
(chapter~11.6) is the most efficient method we have covered.  You
should try rewriting DETECTOR using Simpson's rule to see how it
improves the performance.


\newpage
%%%%%%%%%%%%%%%%%%%%%%%%%%%%%%%%%%%%%%%%%%%%%%%%%%%%%%%%%%%%%%%%%%%%%%%
%                                                                     %
%                              Section 4                              %
%                                                                     %
%%%%%%%%%%%%%%%%%%%%%%%%%%%%%%%%%%%%%%%%%%%%%%%%%%%%%%%%%%%%%%%%%%%%%%%

\section{Fourier Series}

In chapter~10.6 we obtained polynomials which were good approximations
to a function over an interval, where ``good'' meant minimizing the
{\em mean squared separation} between the function and the
approximating polynomials.

While polynomials are the most obvious approximating functions to use
due to the ease with
%
\mar{Difficulties with polynomial approximations}
%
which they can be evaluated, we have seen that finding good
approximating polynomials leads to several serious technical
complications.  The first is that we have to solve systems of
equations to determine the unknown coefficients, a procedure that is
very time-consuming, even for a computer if we are trying to get a
polynomial of, say, degree 30.  Further, if we are trying to make the
approximation over even a moderately-sized interval, since we are
evaluating expressions of the form $x^n$, we get large numbers very
rapidly as $x$ and $n$ get large.  This in turn leads to roundoff
problems in the computer routines.
 
Another aspect of these polynomial approximations that makes them
complicated is that the values of the coefficients change as we change
the degree of the approximating polynomial.  Thus if we determine the
least squares fourth-degree approximation and then decide we want the
fifth-degree approximation instead, all the coefficients have to be
recalculated.  Knowing what the coefficient of $x^3$ was in the
fourth-degree approximation is no help at all in knowing what the
coefficient of $x^3$ will be in the fifth-degree approximation.

There are approximating functions of another kind that avoid such
difficulties.  Moreover, these functions are natural ones to use
%
\mar{Approximating\\ periodic functions}
%
when we are trying to approximate {\em periodic} functions. In such
cases it is reasonable to take the simplest periodic functions---sines
and cosine---and try to combine them to approximate more complicated
periodic functions.  This suggests that we want to look at functions
of the form
\begin{align*}
\phi(x) &= a_0 + a_1 \cos{x}+ a_2\cos{2x} + \cdots+ a_n\cos{nx} \\
 &  \mbox{}\quad +b_1 \sin{x} + b_2 \sin{2x} + \cdots + b_n\sin{nx} \\
 &= a_0 + \sum_{k=1}^n a_k \cos{kx} + b_k \sin{kx}.
\end{align*}

Such a combination is called a {\bf trigonometric polynomial of
  \boldmath degree~$n$}\index{polynomial!trigonometric}.  Note that
any function of this form will in fact be periodic with period $2\pi$.
More generally, if we were interested in approximating a function of
period~$T$, we would want to look at trigonometric polynomials of the
form 
%
\mar{Trigonometric polynomial of degree $n$ and period $T$}
%
\begin{align*}
\phi(x) &= a_0 +a_1 \cos\dfrac{2\pi x}{T}+ a_2 \cos\dfrac{4\pi x}{T} 
        + \cdots+ a_n \cos\dfrac{2n\pi x}{T} \\
 & \mbox{}\quad +b_1 \sin\dfrac{2\pi x}{T}+ b_2 \sin\dfrac{4\pi x}{T} 
        + \cdots+ b_n \sin\dfrac{2n\pi x}{T} \\
 &= a_0+\sum_{k=1}^n a_k \cos\dfrac{2k\pi x}{T} 
        + b_k \sin\dfrac{2k\pi x}{T}.
\end{align*}
You should verify that this does indeed have period~$T$.

To find the coefficients $a_k$ and $b_k$ of the trigonometric
polynomial that best fits a period-$2 \pi$ function $f$ over the
interval $[c, c + 2 \pi]$, we proceed exactly as we did in the
previous section, using the least squares criterion.  That is, for a
given degree $n$, we want to find coefficients $a_0, \ldots , a_n$ and
$b_1, \ldots , b_n$ that minimize the integral
$$
\int_c^{c+2\pi} \left( f(x) -\phi (x)\right)^2\, dx. 
$$ 
In practice, $c$ is usually either 0 or $-\pi$.

The solution turns out to be remarkably compact 
%
\mar{Coefficients are independent of $n$} 
%
and easy to state. One of the key features of the formulas for the
coefficients is that they are independent of each other and of the
particular value of $n$ being used.  Thus, for example, $a_3$ in the
7-th degree approximation has the same value as $a_3$ in the 39-th
degree approximation.  This is a major advantage compared to the
polynomial approximations over intervals that we worked with in
chapter~10.

\smallskip

\begin{center}\label{fourforms}
\framebox[5in]
{\parbox{4.5in}{
\vspace*{6pt} For a function $f$ with period $2 \pi$, its least
squares $n$th degree trigonometric polynomial approximation over a
full period is\vspace{-10pt}
$$
\phi_n(x) = a_0 + \sum_{k=1}^n  a_k \cos{kx}
                  + b_k \sin{kx}, \vspace{-20pt}
$$
where\vspace{-10pt}
\begin{align*}
a_0 &= \dfrac{1}{2\pi} \int_0^{2\pi}f(x)\, dx, \\[3pt]
a_k &= \dfrac{1}{\pi} \int_0^{2\pi}f(x)\cdot \cos{kx}\, dx \quad 
              \mbox{for $k=1, 2, \ldots , n$},\\[3pt]
b_k &= \dfrac{1}{\pi} \int_0^{2\pi}f(x)\cdot \sin{kx}\, dx \quad 
              \mbox{for $k=1, 2, \ldots , n$}.
\end{align*}\vspace{-10pt}}
}
\end{center}

\newpage

The infinite series
$$
a_0 + \sum_{k=1}^{\infty} a_k \cos{kx} + b_k \sin{kx}
$$ 
with the prescribed values for $a_k$ and $b_k$ is called the {\bf
  Fourier series}\index{Fourier series} for~$f$, and the coefficients
$a_k$ and $b_k$ are called the {\bf Fourier
  coefficients}\index{Fourier coefficients} for~$f$.  It turns out
that any continuous function equals its Fourier series in the same
sense we used earlier with Taylor series---for any $x$ in the given
interval, $f(x)$ is the limit as $n\rightarrow\infty$ of the $n$th
degree approximating trigonometric polynomials evaluated at $x$.
The derivation is straightforward, but we shall leave it to the end
of this section so we can look at some examples first.

\aside{Joseph Fourier (1768--1830)\index{Fourier, J.}
  was active in both politics and in mathematics.  He was an advocate
  of the French Revolution, worked as an engineer in Napoleon's army,
  and served as a prefect for a while.  In mathematics he was
  interested in the mathematics of heat conduction and developed the
  series that now bear his name as a tool for investigating problems
  in this area.  His ideas initially met with considerable resistance,
  but eventually became a central tool in mathematics.}

Although our formulas give the values of $a_k$ and $b_k$ in terms of
integrals over $[0, 2\pi]$, periodicity of the integrands implies that
integrations over \emph{any} interval of width $2\pi$ gives the same
values. In practice (as in the first example, immediately below), we
often use $[-\pi, \pi]$ instead of $[0, 2\pi]$.
 
\medskip
\noindent
{\bf Example 1}\quad\label{triangle function} Let's find the
approximating trigonometric polynomials for
\[
f(x) = \begin{cases}
\pi + x & \quad\text{if $-\pi\le x\le 0$}, \\
\pi - x & \quad\text{if $0\le x\le \pi$}.
\end{cases}
\]
Then the graph of $f$
%
\mar{\vspace{.25in}The graph of $f$\\ is ``triangular"}
%
simply consists of two line segments:

\vspace*{1.44in}
\special{psfile=../ps/calc2/sawtoo1.ps}

%\noindent
Now make $f(x)$ periodic over the entire $x$-axis by horizontal
translations: $f(x) = f(x - 2\pi)$. The periodic graph is shown in
gray, above.  We can obtain first Fourier coefficient without any
calculus at all: $a_0=(1/2\pi)\,\pi^2=\pi/2$. It is just the area of
one triangle divided by $2\pi$.  The other coefficients can be
evaluated with integration by parts (chapter~11.3).  We
have
\begin{align*}
a_k &= \dfrac{1}{\pi}\int_{-\pi}^{\pi} f(x) \cos{kx} \, dx \\
    &= \dfrac{1}{\pi}\int_{-\pi}^0 (\pi+x) \cos{kx} \, dx 
         + \dfrac{1}{\pi}\int_0^{\pi}(\pi-x) \cos{kx} \, dx.
   \end{align*}
The first of these integrals can be evaluated as
\begin{align*}
\int_{-\pi}^0 (\pi+x)\cos{kx}\, dx  
  &= \left.(\pi+x)\dfrac{\sin{kx}}{k}\, \right|_{-\pi}^0
       - \int_{-\pi}^0 \dfrac{\sin{kx}}{k}\, dx \\
  &= 0 + \left.\dfrac{\cos{kx}}{k^2}\, \right|_{-\pi}^0 \\
  &= \begin{cases}
\dfrac{2}{k^2} & \quad\text{if $k$ is odd}, \\[6pt]
0  & \quad\text{if $k$ is even}.
     \end{cases}
\end{align*}
Similarly we find
\begin{align*}
\int_0^{\pi}(\pi-x)\cos{kx} \, dx  
  &= \left.(\pi-x)\dfrac{\sin{kx}}{k}\, \right|_0^{\pi}
      + \int_0^{\pi} \dfrac{\sin{kx}}{k}\, dx \\
  &= 0 - \left.\dfrac{\cos{kx}}{k^2}\, \right|_0^{\pi} \\
  &=\begin{cases}
\dfrac{2}{k^2} & \quad\text{if $k$ is odd}, \\[6pt]
0  & \quad\text{if $k$ is even}.
\end{cases}
\end{align*}
   
Combining these two integrals we find
\[
a_k = \begin{cases}
\dfrac{4}{\pi k^2} & \quad\text{if $k$ is odd}, \\[6pt]
\hphantom{0}0  & \quad\text{if $k$ is even}.
\end{cases}
\]
An analogous derivation will show that all the $b_k$ are~0; this is
left to the exercises. We can thus write down the Fourier series
%
\mar{\vspace{.25in}The Fourier series for $f$}
%
for~$f$: 
\[
f(x) = \dfrac{\pi}{2} +\dfrac{4}{\pi}\left(\dfrac{\cos{x}}{1} 
        + \dfrac{\cos{3x}}{9} + \cdots 
        +\dfrac{\cos{(2n+1)x}}{(2n+1)^2} + \cdots \right). 
\]

Let
$$ 
\phi_n(x) = \dfrac{\pi}{2} + \dfrac{4}{\pi} 
     \sum_{k=0}^n \dfrac{\cos{(2k+1)x}}{(2k+1)^2}.
$$
Here are the graphs of $\phi_1(x)$, $\phi_2(x)$, and $\phi_{10}(x)$:

\vspace*{1.5in}
\special{psfile=../ps/calc2/sawtoo2.ps}

\noindent
We see that $\phi_{10}(x)$ already appears to be a very good
approximation to~$f(x)$.  If we look at the maximum separation between
$f(x)$ and $\phi_n(x)$ over $[-\pi, \pi]$ for different values of $n$,
we get the following:
$$
\begin{array}{c|cccccc}
n  & 1&2&10&50&100&1000 \\ [2pt] \hline
	&&&&&& \\ [-8pt]
{\displaystyle \max_{-\pi\le x\le \pi}|f(x)-\phi_n(x)|}
        &.298&.156&.032&.0064&.0032&.00032
\end{array}
$$

Since $\phi_n(x)$ is periodic, 
%
\mar{Approximating a triangular wave-form}
%
if we graph it over a larger interval, we get an approximation to a
{\bf triangular wave-form}\index{wave-form!triangular}.  Here, for
example, is the graph of $\phi_{20}(x)$ over the interval $[-\pi,
  5\pi]$:


\vspace{100pt}
\special{psfile=../ps/calc2/sawtoo3.ps}

\noindent
{\bf Remark 2}\quad In addition to their use in approximating functions,
Fourier series can lead to some interesting, and non-obvious,
mathematical results.  For instance in the preceding example, we have
$f(0)=\pi$\,. On the other hand, we should get the same value if we
set $x=0$ in the Fourier series for $f$.  This leads to the identity
$$
\pi = \dfrac{\pi}{2} + \dfrac{4}{\pi}
      \left( \dfrac{1}{1} + \dfrac{1}{9} + \dfrac{1}{25} 
       + \dfrac{1}{49} + \dfrac{1}{81} + \cdots \right).
$$
With a little rearranging, this can be rewritten as
$$
\dfrac{\pi^2}{8} = 1 + \dfrac{1}{9} + \dfrac{1}{25} 
   + \dfrac{1}{49} + \dfrac{1}{81} + \cdots 
$$
---that is, if we add up the reciprocals of the squares of all the odd
integers, we get $\pi^2/8$!

\medskip

The formulas given on page~\pageref{fourforms} 
%
\mar{The general rule\\ for calculating\\ Fourier series} 
%
for approximating functions with period $2\pi$ extend readily to
approximating periodic functions of any period~$T$.  For instance, if
we wanted to approximate some function $f$ over the interval $[0, T]$,
we have the following formulas.  Check that when $T=2\pi$, these
equations reduce to the earlier ones.  Again, there is nothing special
about the interval $[0, T]$.  If we had wanted to make the
approximation over any other interval of length $T$---for example,
$[-T/2, T/2]$---we simply change the limits of integration to be the
endpoints of that interval.


\begin{center}
\framebox[5in]
{\parbox{4.6in}{
\vspace*{4pt} For a function $f(t)$ with period $T$, its least squares
$n$th degree trigonometric polynomial approximation over a full period
is\vspace{-4pt}
\[
\phi_n(x) =a_0+\sum_{k=1}^n a_k \cos\dfrac{2k\pi x}{T}
     + b_k \sin\dfrac{2k\pi x}{T},\vspace{-10pt}
\]
where\vspace{-10pt}
\begin{align*}
a_0 &= \dfrac{1}{T}\int_0^T f(x)\, dx, \\[8pt]
a_k &= \dfrac{2}{T}\int_0^T f(x)\cdot \cos\dfrac{2k\pi x}{T}\, dx 
            \quad \mbox{for $k=1,2,\ldots ,n$},  \\[8pt]
b_k &= \dfrac{2}{T}\int_0^T f(x)\cdot \sin\dfrac{2k\pi x}{T}\, dx 
            \quad \mbox{for $k=1,2,\ldots ,n$}.
\end{align*}\vspace{-8pt}}
}
\end{center}

\medskip
\noindent
{\bf Example 2}\quad\label{May Fourier} Consider the 
%
\mar{Fourier series for the periodic solutions of May's predator-prey model} 
%
predator--prey model\index{predator--prey model} of May\index{May, R.} that
we examined in chapter~7.3.  Recall that there were two species, the
predator $y$ and the prey $x$, interacting according to the model
\begin{align*}
\mbox{prey:} \quad x' 
   &= .6\,x\left(1-\frac{x}{10}\right) - \frac{.5\,xy}{x+1}, \\
\mbox{predator:} \quad y' 
   &= .1\,y\left(1-\frac{y}{2x}\right).
\end{align*}
We discovered that the populations seemed to move towards periodic
cycles, regardless of the initial conditions (although the phase of
the cycles did depend on the starting values).  In particular, if we
begin with values on this cycle, our solution should be perfectly
periodic with period, it turned out, $T=38.6$ days.  So let's start
with $x=7.75$ and $y=2.38$, values that put us at the peak of the prey
cycle.  Now go to the differential equations and compute the solution
numerically, storing the $x$-values as an array. We can then use these
values to calculate all the integrals needed to find the Fourier
coefficients to approximate the function $x(t)$.  Here are the first
13 terms of the series:
\begin{align*}
x(t)&= 3.7951 + 3.8125\cos\dfrac{2\pi x}{T} + .1514\cos\dfrac{4\pi x}{T}
           + .0326\cos\dfrac{6\pi x}{T} \\[6pt]
 & \mbox{}\quad - .0303\cos\dfrac{8\pi x}{T} 
      - .0609\cos\dfrac{10\pi x}{T}+.0308\cos\dfrac{12\pi x}{T} 
      + \cdots \\[8pt]
 & \mbox{}\quad + 1.1724\sin\dfrac{2\pi x}{T}
      - .0867\sin\dfrac{4\pi x}{T}-.3954\sin\dfrac{6\pi x}{T} \\[8pt]
 & \mbox{}\quad + .0639\sin\dfrac{8\pi x}{T}
      - .0142\sin\dfrac{10\pi x}{T}+.0129\sin\dfrac{12\pi x}{T} 
      + \cdots .
\end{align*}


Let $\phi_3(t)$ be the 7-term trigonometric polynomial whose final
terms involve $\cos(6\pi t/T)$ and $\sin(6\pi t/T)$. Below we graph
$\phi_3(t)$ (dashed line) and $x(t)$ (solid gray line) together.  They
are almost indistinguishable; we let $\phi_3(t)$ run on a little
beyond $x(t)$ so you can see it's there.

\vspace*{170pt}
\special{psfile=../ps/calc2/BMay.ps}

\newpage


\subsubsection{Derivation of the formula for the Fourier coefficients}

The logic behind the derivation is the same as that used in the
previous subsection to find the least squares polynomial
approximations. Fix $n$ and let
$$
\phi(x) = a_0 + \sum_{k=1}^n a_k \cos\dfrac{2k\pi x}{T} 
     + b_k \sin\dfrac{2k\pi x}{T},
$$
where now we want to choose values of the $a_k$ and $b_k$ to minimize
the integral
$$
\int_0^T \left( f(x) -\phi (x)\right)^2\, dx.
$$ 
The value of this integral is thus a function of the undetermined
coefficients $a_0, \ldots , a_n$ and $b_1, \ldots , b_n$.  To find the
coefficients that minimize that value we calculate the partial
derivatives with respect to $a_0$, $a_1$, \ldots as before and set
them equal to 0.  

\enlargethispage*{15pt}
Note that
$$
\dfrac{\partial}{\partial a_m}\phi(x) = \cos\dfrac{2m\pi x}{T}
   \quad \mbox{and} \quad 
\dfrac{\partial}{\partial b_m}\phi(x) = \sin\dfrac{2m\pi x}{T},
$$
so that
\[
\dfrac{\partial}{\partial a_m}
   \int_0^T (f(x) -\phi (x))^2 \, dx 
 = \int_0^T 2(f(x) -\phi (x))
                \left(-\cos\dfrac{2m\pi x}{T}\right)\, dx
\]
and
\[
\dfrac{\partial}{\partial b_m}
\int_0^T (f(x) -\phi (x))^2 \, dx 
 = \int_0^T 2 (f(x) -\phi (x))
                \left(-\sin\dfrac{2m\pi x}{T}\right)\, dx.
\]
The condition 
that all the partial derivatives must be 0 thus leads to the equations
\begin{align*}
\int_0^T 2 (f(x) - \phi(x))(-1) \, dx 
  &= 0, \\[3pt]
\int_0^T 2 (f(x) - \phi(x))
                \left(-\cos\dfrac{2\pi x}{T}\right) \, dx 
  &= 0, \\[6pt]
\int_0^T 2 (f(x) - \phi(x))
                \left(-\cos\dfrac{4\pi x}{T}\right) \, dx 
  &= 0, \\[-5pt]
  &\ \, \vdots \\[-5pt]
\int_0^T 2 (f(x) - \phi(x))
                \left(-\cos\dfrac{2n\pi x}{T}\right) \, dx 
  &= 0, 
\end{align*}
and
\begin{align*}
\int_0^T 2 (f(x) - \phi(x))
                \left(-\sin\dfrac{2\pi x}{T}\right) \, dx
  &= 0, \\[8pt]
\int_0^T 2 (f(x) - \phi(x))
                \left(-\sin\dfrac{4\pi x}{T}\right) \, dx
  &= 0, \\[-3pt]
  &\ \, \vdots \\[-3pt]
\int_0^T 2 (f(x) - \phi(x))
                \left(-\sin\dfrac{2n\pi x}{T}\right) \, dx
  &= 0.
\end{align*}
These equations can be rewritten as
\begin{align*}
\int_0^T f(x) \, dx
  &= \int_0^T \phi(x)\, dx,  \\[10pt]
\int_0^T f(x)\cos\dfrac{2\pi x}{T} \, dx
  &= \int_0^T \phi(x)\cos\dfrac{2\pi x}{T}\, dx,  \\[10pt]
\int_0^T f(x)\cos\dfrac{4\pi x}{T} \, dx
  &= \int_0^T \phi(x)\cos\dfrac{4\pi x}{T}\, dx,  \\[-3pt]
  &\ \, \vdots  \\[-3pt]
\int_0^T f(x)\cos\dfrac{2n\pi x}{T} \, dx
  &= \int_0^T \phi(x)\cos\dfrac{2n\pi x}{T}\, dx,
\end{align*}
and
\begin{align*}
\int_0^T f(x)\sin\dfrac{2\pi x}{T} \, dx
  &= \int_0^T \phi(x)\sin\dfrac{2\pi x}{T}\, dx,  \\[10pt]
\int_0^T f(x)\sin\dfrac{4\pi x}{T} \, dx
  &= \int_0^T \phi(x)\sin\dfrac{4\pi x}{T}\, dx,  \\[-3pt]
  &\ \, \vdots  \\[-3pt]
\int_0^T f(x)\sin\dfrac{2n\pi x}{T} \, dx
  &= \int_0^T \phi(x)\sin\dfrac{2n\pi x}{T} \, dx,
\end{align*}

The Fourier coefficients $a_k$ and $b_k$ that we seek appear
in~$\phi$, and we shall obtain them by calculating the integrals on
the right (the ones involving~$\phi$) in the equations above.  Since 
\[
\phi(x) = a_0 + \sum_{k=1}^n a_k \cos\dfrac{2k\pi x}{T} 
     + \sum_{k=1}^n b_k \sin\dfrac{2k\pi x}{T},
\]
we have (for each $m = 0, 1, \ldots , n$)
\begin{align*}
\phi(x)\cos\dfrac{2m\pi x}{T}  
   &= a_0 \cos\dfrac{2m\pi x}{T}  
     + \sum_{k=1}^n a_k \cos\dfrac{2k\pi x}{T} \cos\dfrac{2m\pi x}{T} \\  
   &\mbox{}\quad 
     + \sum_{k=1}^n b_k \sin\dfrac{2k\pi x}{T} \cos\dfrac{2m\pi x}{T},  \\
\intertext{and (for each $m = 1, \ldots , n$)}
\phi(x)\sin\dfrac{2m\pi x}{T}  
   &= 
      \sum_{k=1}^n a_k \cos\dfrac{2k\pi x}{T} \sin\dfrac{2m\pi x}{T}
     + \sum_{k=1}^n b_k \sin\dfrac{2k\pi x}{T} \sin\dfrac{2m\pi x}{T}.  
\end{align*}


The integrals of the expressions on the left 
%
\mar{Integrating products\\ of sines and cosines}
%
therefore reduce to sums of integrals of various products of sines and
cosines.  In each sum, only \emph{one} term yields a nonzero integral.
All the values are given below.  The formulas are left for you to derive
in the exercises, using integration formulas from the exercises in
chapter~11.3.  For integers $k$ and $m$, we have
\begin{align*}
\int_0^T \sin\dfrac{2k\pi x}{T} \cos\dfrac{2m\pi x}{T} \, dx 
  &=  0 \quad \mbox{ for all $k$ and $m$}; 
  \\[6pt]
\int_0^T \sin\dfrac{2k\pi x}{T} \sin\dfrac{2m\pi x}{T} \, dx 
  &= \begin{cases}
	T/2 & \text{if $k=m$}, \\
	0   & \text{otherwise};
     \end{cases}
 \\[8pt]
\int_0^T \cos\dfrac{2k\pi x}{T} \cos\dfrac{2m\pi x}{T} \, dx 
  &= \begin{cases}
	T/2  & \text{if $k=m \ne 0$}, \\
	T & \text{if $k=m=0$}, \\
	0    & \text{otherwise}.
     \end{cases}
\end{align*}
For $m = 0$ we have $\cos\dfrac{2m\pi x}{T} = \cos{0} = 1$, so
\[
\int_0^T \phi(x)\cos\dfrac{2m\pi x}{T} \, dx
   = \int_0^T a_0  \ dx
   = a_0 \cdot T;
\]
it follows that 
\[
a_0 = \dfrac{1}{T} \int_0^T \phi(x) \, dx.
\]
For each $m = 1, 2, \ldots , n$, we have, first of all,
\[
\int_0^T \phi(x)\cos\dfrac{2m\pi x}{T} \, dx 
   = \int_0^T a_m \cos^2\dfrac{2m\pi x}{T} \, dx 
   = a_m \cdot T/2,
\]
from which it follows that
\[
a_m = \dfrac{2}{T} \int_0^T \phi(x)\cos\dfrac{2m\pi x}{T} \, dx.
\]
Second, we have
\[
\int_0^T \phi(x)\sin\dfrac{2m\pi x}{T} \, dx 
   = \int_0^T b_m \sin^2\dfrac{2m\pi x}{T} \, dx 
   = b_m \cdot T/2,
\]
so
\[
b_m = \dfrac{2}{T} \int_0^T \phi(x)\sin\dfrac{2m\pi x}{T} \, dx.
\]
The derivation is complete.

\subsection{Exercises}

\vspace{-10pt}

\num Use the formulas on page~\pageref{`trig'} in chapter~11.3 to
derive the following equalities; $k$ and $m$ are integers.

\medskip

\sub $\dint_0^T \sin\dfrac{2k\pi x}{T} \cos\dfrac{2m\pi x}{T} \, dx 
  =  0$ for all $k$ and $m$.

\medskip

\sub $\dint_0^T \sin\dfrac{2k\pi x}{T} \sin\dfrac{2m\pi x}{T} \, dx 
  = \begin{cases}
	T/2 & \text{if $k=m$}, \\
	0   & \text{otherwise}.
     \end{cases}$

\medskip

\sub  $\dint_0^T \cos\dfrac{2k\pi x}{T} \cos\dfrac{2m\pi x}{T} \, dx 
  = \begin{cases}
	T/2  & \text{if $k=m \ne 0$}, \\
	T & \text{if $k=m=0$}, \\
	0    & \text{otherwise}.
     \end{cases}$

\num Show that in the Fourier series for the triangular function
discussed in the text (example ~1, page~\pageref{triangle function}),
all the coefficients of the sine terms really are~0.

%\newpage

\num Find the Fourier series for the following functions over the
interval $[-\pi,\pi]$:

\sub $f(x)=x$. \hfill [Ans. ${\displaystyle
    2\sum_{k=1}^{\infty}\dfrac{(-1)^{n-1}\sin nx}{n}}$]

\sub $f(x)= \pi^2-x^2$.

\sub $f(x) = \begin{cases}
0 & \mbox{if $-\pi\le x\le 0$}, \\ 
x^2  & \mbox{if $0\le x\le \pi$}.
\end{cases}$

\num In May's predator--prey model, find the first seven terms of the
Fourier series for the predator species,~$y(t)$.  Use $T = 38.6$ days
and the initial conditions $x = 7.75$ and $y = 2.38$, and in example~2
in the text (page~\pageref{May Fourier}).\index{May model}


\cleardoublepage


